
\renewcommand{\thefigure}{S-\arabic{figure}}
\setcounter{figure}{0}
\renewcommand{\thetable}{S-\arabic{table}}
\setcounter{table}{0}
\renewcommand{\thesection}{S-\arabic{section}}
\setcounter{section}{0}
\makeatletter
\@ifundefined{theHfigure}
{\newcommand{\theHfigure}{S-\arabic{figure}}}
{\renewcommand{\theHfigure}{S-\arabic{figure}}}
\@ifundefined{theHtable}
{\newcommand{\theHtable}{S-\arabic{table}}}
{\renewcommand{\theHtable}{S-\arabic{table}}}
\@ifundefined{theHsection}
{\newcommand{\theHsection}{S-\arabic{section}}}
{\renewcommand{\theHsection}{S-\arabic{section}}}
\makeatother
% \renewcommand{\theequation}{S\arabic{equation}}
% \setcounter{equation}{0}

\FloatBarrier
\begin{center}
    {\LARGE\bfseries Supplementary Information for

        \textit{DPA4: Pushing the Accuracy--Cost Frontier of Interatomic
            Potentials with EMFA SO(2) Convolution} \par}
    \vspace{1em}
\end{center}

\section{Mathematical notation and equivariant operators}
\label{sec:sm-model-formulation}
\label{sec:sm-methods}

This section collects the mathematical details that supplement the architecture description in Section~\ref{sec:methods}: the formal SO(3) representation notation, the edge-local frame and SO(2) decomposition, the closed-form cutoff envelope coefficients, the truncated SO(2) layout and rescaled lift, the SO(2)-equivariant operator algebra, and the cutoff-consistent first-layer bias correction used by the ablations in later supplementary sections.

\subsection{SO(3) representation: notation}
\label{sec:sm-representation}

For each integer $l\ge 0$, let $V_l\cong\mathbb{R}^{2l+1}$ be the real $(2l+1)$-dimensional irreducible representation of SO(3), spanned by the real spherical harmonics $\{Y_l^m\}_{m=-l}^{l}$.
The action of a rotation
$Q\in\mathrm{SO}(3)$ on $V_l$ is given by the real Wigner $D$-matrix
$D^l(Q)\in\mathrm{O}(2l+1)$, the orthogonal $(2l+1)\times(2l+1)$ matrix
that describes how the real spherical harmonics transform under rotation,
\begin{equation}
    Y_l^m(Q^{-1}\widehat{\mathbf{r}}) =\sum_{m'=-l}^{l}D^l(Q)_{m,m'}\,Y_l^{m'}(\widehat{\mathbf{r}}).
    \label{eq:sm-wigner-D}
\end{equation}
The map $D^l:\mathrm{SO}(3)\to\mathrm{O}(2l+1)$ is a continuous group homomorphism (i.e.\ $D^l(QQ')=D^l(Q)D^l(Q')$), and the $V_l$ are \emph{irreducible}: $V_l$ admits no proper SO(3)-invariant subspace.
The trivial case $l=0$ is the invariant scalar representation with $D^0(Q)\equiv 1$.

The node-feature space $V_{\le L}\otimes\mathbb{R}^{C}=\bigoplus_{l=0}^{L}
    V_l\otimes\mathbb{R}^{C}$ of main-text Section~\ref{sec:methods} therefore
carries the real block-diagonal action
\begin{equation}
    D(Q)=\bigoplus_{l=0}^{L}D^l(Q), \label{eq:sm-direct-sum} \end{equation} acting independently on each degree-$l$ block and trivially on the channel factor $\mathbb{R}^{C}$.
Basis coefficients are packed by the linear index
\begin{equation}
    \iota(l,m)=l^2+l+m,\qquad
    m=-l,\ldots,l,\quad l=0,\ldots,L,
    \label{eq:sm-pack-index}
\end{equation}
used throughout this Supplement and in main-text Eq.~\eqref{eq:dpa4-gie};
under this packing, the matrix $D(Q)$ is block-diagonal with the
$(2l+1)\times(2l+1)$ block $D^l(Q)$ occupying rows and columns
$\iota(l,-l),\ldots,\iota(l,+l)$.
A map $f:V_{\le L}\otimes\mathbb{R}^{C}\to V_{\le L}\otimes\mathbb{R}^{C'}$ is \emph{SO(3)-equivariant} if $f\circ D(Q)=D(Q)\circ f$ for every $Q\in\mathrm{SO}(3)$, and \emph{SO(3)-invariant} if it lands in the $l=0$ block, on which $D^0\equiv 1$.

\subsection{Edge-local frame and SO(2) decomposition}
\label{sec:sm-local-frame}

\paragraph{Goal.}
For each directed edge $(i,j)$, DPA4 chooses a rotation $R_{ij}$ satisfying
\begin{equation}
    R_{ij}\widehat{\mathbf{r}}_{ij}=\mathbf{e}_z=(0,0,1)^{\top}.
    \label{eq:sm-local-frame-goal}
\end{equation}
This converts an SO(3) problem into an SO(2) problem in the edge-local frame: after the bond direction is aligned with $\mathbf{e}_z$, the residual symmetry is the subgroup of rotations about $\mathbf{e}_z$.
The in-plane basis is a gauge choice.
The equivariant message does not depend on this choice because the local operator commutes with the residual SO(2) action.

\paragraph{Two quaternion charts.}
Let $\widehat{\mathbf{r}}=(x,y,z)\in \Sph$.
DPA4 uses two smooth unit-quaternion
charts,
\begin{equation}
    \mathbf{q}^{+}(\widehat{\mathbf{r}})
    =
    \frac{(1+z,\;y,\;-x,\;0)}{\sqrt{2(1+z)}},
    \qquad
    \mathbf{q}^{-}(\widehat{\mathbf{r}})
    =
    \frac{(-x,\;0,\;1-z,\;y)}{\sqrt{2(1-z)}}.
    \label{eq:sm-quaternion-charts}
\end{equation}
The first chart is regular away from the south pole and the second is regular away from the north pole.
In the overlap, the sign of $\mathbf{q}^{-}$ is chosen
so that $\langle\mathbf{q}^{+},\mathbf{q}^{-}\rangle\ge 0$, and the blended
quaternion is
\begin{equation}
    \mathbf{q}_{ij}
    =
    \frac{\lambda\mathbf{q}^{+}(\widehat{\mathbf{r}}_{ij})
        +(1-\lambda)\mathbf{q}^{-}(\widehat{\mathbf{r}}_{ij})}
    {\left\|\lambda\mathbf{q}^{+}(\widehat{\mathbf{r}}_{ij})
        +(1-\lambda)\mathbf{q}^{-}(\widehat{\mathbf{r}}_{ij})\right\|},
    \qquad
    \lambda=\frac{1+z}{2}.
    \label{eq:sm-quaternion-blend}
\end{equation}
The denominator is bounded away from zero after shortest-arc sign alignment in the chosen chart overlap, so the resulting gauge is smooth on the chart used by the implementation.
The Wigner-D matrix in the edge-local gauge is denoted $D_{ij}=D(R(\mathbf{q}_{ij}))$.
The underlying unit-quaternion rotation matrix is
\begin{equation}
    R(\mathbf{q})=
    \begin{pmatrix}
        1-2(q_y^2+q_z^2) & 2(q_xq_y-q_wq_z) & 2(q_xq_z+q_wq_y) \\
        2(q_xq_y+q_wq_z) & 1-2(q_x^2+q_z^2) & 2(q_yq_z-q_wq_x) \\
        2(q_xq_z-q_wq_y) & 2(q_yq_z+q_wq_x) & 1-2(q_x^2+q_y^2)
    \end{pmatrix}.
    \label{eq:sm-quaternion-rotation}
\end{equation}
During training, an additional random roll about the local $z$ axis may be composed with $R_{ij}$.
This roll is a gauge augmentation: it changes the in-plane basis of the local gauge but leaves the bond direction fixed.
Because the SO(2) stack is gauge equivariant, the lifted global message is unchanged by this roll apart from the prescribed SO(3) transformation law.

\paragraph{SO(2) decomposition.}
Let $\{\mathbf{e}^{(l)}_m\}_{m=-l}^{l}$ denote the real-spherical-harmonic basis of $V_l$.
Under restriction to the SO(2) subgroup of rotations about
the local $z$ axis, $V_l$ decomposes as
\begin{equation}
    V_l|_{\mathrm{SO}(2)} = \mathrm{span}_{\mathbb{R}}\{\mathbf{e}^{(l)}_0\} \;\oplus\; \bigoplus_{m=1}^{l} \mathrm{span}_{\mathbb{R}}\{\mathbf{e}^{(l)}_{-m},\,\mathbf{e}^{(l)}_{+m}\}, \label{eq:sm-so2-decomposition} \end{equation} where $\mathrm{span}_{\mathbb{R}}\{\cdot\}$ denotes the real linear span of the indicated basis vectors.
The $m=0$ summand is a one-dimensional trivial SO(2) sub-representation, and each $m=1,\ldots,l$ summand is a two-dimensional real SO(2) sub-representation on which a $z$-axis rotation by angle $\theta$ acts as planar rotation by angle $m\theta$, equivalent to multiplication by the complex phase $e^{im\theta}$.
This is the algebraic reason DPA4 can use SO(2)-equivariant local operators instead of Clebsch--Gordan tensor products.

\subsection{Closed-form cutoff coefficients}
\label{sec:sm-radial}

The four boundary conditions
$s_p(r_{\mathrm{c}})=s_p'(r_{\mathrm{c}})=s_p''(r_{\mathrm{c}})=s_p'''(r_{\mathrm{c}})=0$
of the cutoff envelope $s_p(r)$ in main-text
Eq.~\eqref{eq:dpa4-envelope} uniquely fix the coefficients to
\begin{align}
    a_p & =-\frac{(p+1)(p+2)(p+3)}{6}, &
    b_p & =\frac{p(p+2)(p+3)}{2},        \\
    c_p & =-\frac{p(p+1)(p+3)}{2},     &
    d_p & =\frac{p(p+1)(p+2)}{6}.
    \label{eq:sm-cutoff-coefficients}
\end{align}
For the two values used in DPA4 ($s_5$ for edge weighting and the smooth
degree, $s_7$ inside the radial basis of main-text
Eq.~\eqref{eq:dpa4-radial-basis}), the explicit polynomials are
\begin{align}
    s_5(r) & =1-56x^5+140x^6-120x^7+35x^8,       \\
    s_7(r) & =1-120x^7+540x^8-1080x^9+840x^{10},
\end{align}
with $x=r/r_{\mathrm{c}}$.

\subsection{Truncated local layout and rescaled lift}
\label{sec:sm-truncation}

Inside the edge-local frame, DPA4 retains only
\begin{equation}
    \mathcal{I}_M=
    \{(l,m):0\le l\le L,\ |m|\le \min(l,M)\},
    \qquad
    D_M=|\mathcal{I}_M|.
    \label{eq:sm-trunc-index}
\end{equation}
Let $P_M:\mathbb{R}^{(L+1)^2}\to\mathbb{R}^{D_M}$ be the orthogonal projection onto this reduced layout.
The truncated edge rotation is
\begin{equation}
    D_{ij}^{\le M}=P_M D_{ij}.
    \label{eq:sm-truncated-rotation}
\end{equation}
Because $P_M^{\top}P_M$ is not the identity when $M<L$, the round trip loses degree-block norm.
DPA4 applies the diagonal lift compensation
\begin{equation}
    (\Xi_M)_{\iota(l,m),\iota(l,m)}
    =
    \kappa_l
    =
    \sqrt{\frac{2l+1}{2\min(l,M)+1}}.
    \label{eq:sm-rescale}
\end{equation}
Thus $\Xi_M= \operatorname{diag}(\kappa_0 I_1,\kappa_1 I_3,\ldots,\kappa_L I_{2L+1})$.

\subsection{Equivariant operator algebra}
\label{sec:sm-operators}

\paragraph{Classification by Schur's lemma.}
The linear operators used by DPA4 follow from the classification of equivariant maps on the relevant representation spaces.
In the global
SO(3) frame, Schur's lemma forbids mixing distinct degrees:
\begin{equation}
    \mathrm{Hom}_{\mathrm{SO}(3)}
    (V_{\le L}\otimes\mathbb{R}^{C},V_{\le L}\otimes\mathbb{R}^{C'})
    \cong
    \bigoplus_{l=0}^{L}\mathbb{R}^{C\times C'}.
    \label{eq:sm-schur-so3}
\end{equation}
In the edge-local SO(2) frame, the $m=0$ lines are trivial representations
and each $|m|>0$ pair is complex type, giving
\begin{equation}
    \mathrm{Hom}_{\mathrm{SO}(2)}
    (V_{\le L}\otimes\mathbb{R}^{C},V_{\le L}\otimes\mathbb{R}^{C'})
    \cong
    \bigoplus_{l,l'}\mathbb{R}^{C\times C'}
    \oplus
    \bigoplus_{m=1}^{M}\bigoplus_{l,l'\ge m}\mathbb{C}^{C\times C'}.
    \label{eq:sm-schur-so2}
\end{equation}
Equation~\eqref{eq:sm-schur-so2} is the formal reason the local-frame operator may mix degrees $l,l'$ while preserving each $|m|$ stratum.

\paragraph{Global SO(3) linear maps.}
By Eq.~\eqref{eq:sm-schur-so3}, an SO(3)-equivariant \emph{degree-wise} map
has the explicit form
\begin{equation}
    (L^{\mathrm{deg}}_{\Theta}\mathbf{h})_{\iota(l,m),c'}
    =
    \sum_{c=1}^{C}
    W^{(l)}_{c,c'}\mathbf{h}_{\iota(l,m),c}, \label{eq:sm-degree-linear} \end{equation} with one learnable channel matrix $W^{(l)}\in\mathbb{R}^{C\times C'}$ per degree $l$.
Here ``degree-wise'' means \emph{per-$l$ but identical across all $m\in\{-l,\ldots,l\}$ within a degree}: the same $W^{(l)}$ is applied to every $m$-slice of the degree-$l$ block, with different matrices allowed for different $l$.
All ``degree-wise channel-mixing maps'' in main-text Section~\ref{sec:methods} (including $L^{\mathrm{pre}}_{\mathrm{deg}}$, $L^{\mathrm{post}}_{\mathrm{deg}}$, $L^{\mathrm{rad}}_{\mathrm{lift}}$, $L^{\mathrm{ch}}_{\mathrm{in}}$ and $L^{\mathrm{ch}}_{\mathrm{out}}$) are instances of this form.

\paragraph{Edge-local SO(2) linear maps.}
In the edge-local frame, the $m=0$ lines are trivial SO(2) representations, whereas each $|m|>0$ pair is complex type.
The most general SO(2)-equivariant
linear map on the reduced layout is therefore
\begin{align}
    (L_{\Theta}^{\mathrm{SO2}}\mathbf{x})_{(l,0),c'}
     & =
    \sum_{l'=0}^{L}\sum_{c=1}^{C}
    A_{c,c'}^{(l,l',0)}\mathbf{x}_{(l',0),c}
    +b_{0,c'}\delta_{l,0},
    \label{eq:sm-so2-m0} \\
    \begin{pmatrix}
        (L_{\Theta}^{\mathrm{SO2}}\mathbf{x})_{(l,-m),c'}\\
        (L_{\Theta}^{\mathrm{SO2}}\mathbf{x})_{(l,+m),c'}
    \end{pmatrix}
     & =
    \sum_{l'\ge m}\sum_{c=1}^{C}
    \begin{pmatrix}
        U_{c,c'}^{(l,l',m)} & -V_{c,c'}^{(l,l',m)} \\
        V_{c,c'}^{(l,l',m)} & U_{c,c'}^{(l,l',m)}
    \end{pmatrix}
    \begin{pmatrix}
        \mathbf{x}_{(l',-m),c} \\
        \mathbf{x}_{(l',+m),c}
    \end{pmatrix}.
    \label{eq:sm-so2-mn}
\end{align}

\paragraph{Equivariant RMS normalization.}
For $\mathbf{h}\in V_{\le L}\otimes\mathbb{R}^{C}$, define
\begin{equation}
    \sigma^2(\mathbf{h})
    =
    \sum_{l=0}^{L}\sum_{m=-l}^{l}\sum_{c=1}^{C}
    \frac{(\mathbf{h}_{\iota(l,m),c}-\delta_{l,0}\bar{h})^2}
    {(2l+1)(L+1)C},
    \qquad
    \bar{h}=C^{-1}\sum_c \mathbf{h}_{\iota(0,0),c}.
    \label{eq:sm-rms-stat}
\end{equation}
The normalization
\begin{equation}
    (N_{\gamma,\beta}\mathbf{h})_{\iota(l,m),c}
    =
    \gamma_l
    \frac{\mathbf{h}_{\iota(l,m),c}-\delta_{l,0}\bar{h}}
    {\sqrt{\sigma^2(\mathbf{h})+\varepsilon}}
    +\delta_{l,0}\beta_c
    \label{eq:sm-rmsnorm}
\end{equation}
commutes with SO(3), because the numerator is equivariant, the denominator is
invariant under the orthogonal representation $D^l$, and $\gamma_l$ is constant
within each degree block.

\paragraph{Scalar-gated nonlinearity.}
For a smooth scalar nonlinearity $\psi$ and degree-wise gate matrices
$G^{(l)}$, define
\begin{equation}
    (\Gamma_{\psi,G}\mathbf{h})_{\iota(l,m),c}
    =
    \begin{cases}
        \psi(\mathbf{h}_{\iota(0,0),c}),                  & l=0,    \\[2pt]
        \mathbf{h}_{\iota(l,m),c}
        \sigma\!\left(\sum_{c'}
        G^{(l)}_{c',c} \mathbf{h}_{\iota(0,0),c'}\right), & l\ge 1.
    \end{cases}
    \label{eq:sm-gated}
\end{equation}
The gate is a scalar function of the invariant $l=0$ slice, so the operation is SO(3)-equivariant.

\paragraph{Scalar-keyed mixtures.}
If $\mathbf{x}^{(1)},\ldots,\mathbf{x}^{(S)}$ are equivariant tensors of the
same type and $\pi$ is an invariant scalar projection, then
\begin{equation}
    \mathcal{A}(\mathbf{x}^{(1)},\ldots,\mathbf{x}^{(S)};\mathbf{q})
    =
    \sum_{s=1}^{S}\alpha_s\mathbf{x}^{(s)},\qquad
    \alpha_s=
    \frac{\exp\langle\mathbf{q},\pi(\mathbf{x}^{(s)})\rangle}
    {\sum_{s'}\exp\langle\mathbf{q},\pi(\mathbf{x}^{(s')})\rangle}
    \label{eq:sm-scalar-keyed-mixture}
\end{equation}
is equivariant, because the weights are invariant scalars.

\subsection{Bias consistency at the smooth cutoff}
\label{sec:sm-so2conv}

If the first SO(2) linear map of the EMFA SO(2) convolution (main-text Sec.~\ref{sec:so2}; Eqs.~\eqref{eq:sm-so2-m0}, \eqref{eq:sm-so2-mn}) includes an additive $l=0$ bias, that bias must vanish with the same smooth envelope as the rest of the edge message.
Otherwise an edge whose radial contribution has gone to zero could still carry a constant scalar offset.
DPA4 therefore treats the first-layer scalar bias as part of the edge-conditioned response.
If $b_{0,f,c}$ is the scalar bias for focus $f$
and channel $c$, the net contribution to the local $(l,m)=(0,0)$ slot is
adjusted to
\begin{equation}
    b_{0,f,c}\,\widetilde{\rho}_{ij,0,c}\,s_5(r_{ij}),
    \label{eq:sm-bias-consistency}
\end{equation}
which is $C^3$ and vanishes at the cutoff.
This correction is only needed at the first SO(2) layer because subsequent layers operate on already modulated local features.

\subsection{Numerical equivariance of truncated local layouts}
\label{sec:sm-truncated-equivariance}

Table~\ref{tab:sm_s2_truncated_equivariance} extends the full-coefficient comparison of main-text Table~\ref{tab:s2_full_equivariance} to the $m$-truncated local layout used inside the EMFA SO(2) convolution, evaluating the maximum equivariance error of product-grid rules and Lebedev quadrature under random rotations about the local $z$ axis.

\begin{table}[p]
    \centering
    \scriptsize
    \setlength{\tabcolsep}{3pt}
    \caption{
        $m$-truncated \texorpdfstring{$\Sph$}{S2} activation equivariance under random local
        $z$-axis rotations.
        These cases test the reduced local layout used in SO(2) convolution.
        Product-grid rules are $(R_{\phi},R_{\theta})$ with the total number of grid points $R_{\phi}R_{\theta}$ given alongside; Lebedev rules are reported by their algebraic order of accuracy $p$ and the corresponding number of points.
        Errors are maximum absolute deviations between the two equivariance paths.
    }
    \label{tab:sm_s2_truncated_equivariance}
    \begin{tabular*}{\linewidth}{@{\extracolsep{\fill}}cccccccccc@{}}
        \toprule
        \textbf{$M$} &
        \textbf{$L$} &
        \multicolumn{4}{c}{\textbf{Product grid}} &
        \multicolumn{4}{c}{\textbf{Lebedev quadrature}} \\
        \cmidrule(lr){3-6}\cmidrule(lr){7-10}
        & & \textbf{Rule} & \textbf{\#\,pts} & \textbf{fp64 error} & \textbf{fp32 error}
        & \textbf{$p$} & \textbf{\#\,pts} & \textbf{fp64 error} & \textbf{fp32 error} \\
        \midrule
        1 & 2 & $6\times8$   & 48  & $2.36\times10^{-7}$ & $3.58\times10^{-7}$ & 7  & 26  & $2.31\times10^{-14}$ & $2.38\times10^{-7}$ \\
        1 & 3 & $6\times12$  & 72  & $1.22\times10^{-7}$ & $5.96\times10^{-7}$ & 9  & 38  & $3.55\times10^{-14}$ & $2.98\times10^{-7}$ \\
        1 & 4 & $6\times14$  & 84  & $1.12\times10^{-6}$ & $9.54\times10^{-7}$ & 13 & 74  & $1.04\times10^{-13}$ & $9.54\times10^{-7}$ \\
        1 & 5 & $6\times18$  & 108 & $1.10\times10^{-7}$ & $1.43\times10^{-6}$ & 15 & 86  & $9.34\times10^{-14}$ & $7.15\times10^{-7}$ \\
        1 & 6 & $6\times20$  & 120 & $7.64\times10^{-7}$ & $1.91\times10^{-6}$ & 19 & 146 & $8.56\times10^{-14}$ & $2.15\times10^{-6}$ \\
        1 & 7 & $6\times24$  & 144 & $2.17\times10^{-7}$ & $1.91\times10^{-6}$ & 21 & 170 & $2.08\times10^{-13}$ & $3.34\times10^{-6}$ \\
        \midrule
        2 & 2 & $8\times8$   & 64  & $4.01\times10^{-7}$ & $8.34\times10^{-7}$ & 7  & 26  & $1.50\times10^{-14}$ & $2.38\times10^{-7}$ \\
        2 & 3 & $8\times12$  & 96  & $5.99\times10^{-7}$ & $8.34\times10^{-7}$ & 9  & 38  & $5.71\times10^{-14}$ & $3.58\times10^{-7}$ \\
        2 & 4 & $8\times14$  & 112 & $6.02\times10^{-7}$ & $1.67\times10^{-6}$ & 13 & 74  & $9.15\times10^{-14}$ & $5.96\times10^{-7}$ \\
        2 & 5 & $8\times18$  & 144 & $1.19\times10^{-6}$ & $1.55\times10^{-6}$ & 15 & 86  & $7.83\times10^{-14}$ & $4.77\times10^{-7}$ \\
        2 & 6 & $8\times20$  & 160 & $1.33\times10^{-6}$ & $2.15\times10^{-6}$ & 19 & 146 & $1.29\times10^{-13}$ & $9.54\times10^{-7}$ \\
        2 & 7 & $8\times24$  & 192 & $1.41\times10^{-6}$ & $2.62\times10^{-6}$ & 21 & 170 & $1.56\times10^{-13}$ & $1.43\times10^{-6}$ \\
        \bottomrule
    \end{tabular*}
\end{table}

\subsection{Completeness of the Wigner bilinear network}
\label{sec:sm-wbn-complete}

This subsection proves the completeness statement invoked in main-text Section~\ref{sec:s2-method}: networks stacking Wigner bilinear blocks with the \emph{full} frame range $|k|\le l$, interleaved with equivariant linear maps, approximate every continuous SO(3)-equivariant map between band-limited feature spaces, and realize polynomial equivariant maps exactly.
The block analyzed here is the exact continuous model underlying the deployed FFN: it retains all frames, forms products over all pairs of left and right channels, and evaluates the Haar integral exactly.
The deployed block (main-text Eqs.~\eqref{eq:dpa4-wbb-lift}--\eqref{eq:dpa4-ffn}) is its economical specialization -- channel-diagonal products, frames truncated to $|k|\le k_{\max}$, and an exact quadrature in place of the integral -- and the role of the truncation is discussed in Remark~\ref{rem:sm-wbn-truncation}.

Throughout, $g$ denotes a rotation, $dg$ the normalized Haar measure on $SO(3)$, and $D^{(l)}_{mk}(g)$ the \emph{complex} Wigner matrices, which satisfy the orthogonality relation
\begin{equation}
    (2l+1)\int_{SO(3)}\overline{D^{(l)}_{mk}(g)}\,D^{(l')}_{m'k'}(g)\,dg
    =\delta_{ll'}\delta_{mm'}\delta_{kk'}
    \label{eq:sm-wigner-orth}
\end{equation}
and the product identity~\cite{varshalovich1988quantum}
\begin{equation}
    D^{(l_1)}_{m_1k_1}(g)\,D^{(l_2)}_{m_2k_2}(g)
    =\sum_{L=|l_1-l_2|}^{l_1+l_2}
    \langle l_1 m_1\, l_2 m_2|L M\rangle\,
    \langle l_1 k_1\, l_2 k_2|L \kappa\rangle\,
    D^{(L)}_{M\kappa}(g),
    \label{eq:sm-wigner-product}
\end{equation}
with $M=m_1+m_2$, $\kappa=k_1+k_2$, and $\langle\cdot\,\cdot|\cdot\rangle$ the Clebsch--Gordan coefficients.
The associated coupling map $C^{L}_{l_1l_2}:V_{l_1}\otimes V_{l_2}\to V_L$,
\begin{equation}
    \bigl(C^{L}_{l_1l_2}(x,y)\bigr)_M
    =\sum_{m_1,m_2}\langle l_1 m_1\, l_2 m_2|L M\rangle\,x_{m_1}y_{m_2},
    \label{eq:sm-cg-map}
\end{equation}
is a nonzero SO(3)-intertwiner for every admissible triple $|l_1-l_2|\le L\le l_1+l_2$.
The deployed network uses the real Wigner basis of Section~\ref{sec:sm-representation}; the real and complex conventions are related by an invertible equivariant change of coordinates, which is absorbed into the learnable channel maps, so it suffices to argue in the complex convention.

\paragraph{The full-frame block.}
Let $\mathcal{H}=V_{\le L_{\max}}\otimes\mathbb{C}^{C}$.
A \emph{full-frame Wigner bilinear block} $\mathcal{B}:\mathcal{H}\to\mathcal{H}$ maps a feature $x^{(l)}_{m,c}$ to
\begin{align}
    u^{\mathrm{L},(l)}_{m,k,a}
    &=\sum_{c}A^{\mathrm{L},(l)}_{k,a,c}\,x^{(l)}_{m,c},
    \qquad
    u^{\mathrm{R},(l)}_{m,k,b}
    =\sum_{c}A^{\mathrm{R},(l)}_{k,b,c}\,x^{(l)}_{m,c},
    \qquad |k|\le l,
    \label{eq:sm-wbb-lin}\\
    f^{\mathrm{L}}_{a}(g)
    &=\sum_{l,m,k}u^{\mathrm{L},(l)}_{m,k,a}D^{(l)}_{mk}(g),
    \qquad
    f^{\mathrm{R}}_{b}(g)
    =\sum_{l,m,k}u^{\mathrm{R},(l)}_{m,k,b}D^{(l)}_{mk}(g),
    \label{eq:sm-wbb-synth}\\
    h_{ab}(g)
    &=f^{\mathrm{L}}_{a}(g)\,f^{\mathrm{R}}_{b}(g),
    \label{eq:sm-wbb-prod}\\
    z^{(l)}_{m,k,ab}
    &=(2l+1)\int_{SO(3)}\overline{D^{(l)}_{mk}(g)}\,h_{ab}(g)\,dg,
    \label{eq:sm-wbb-analysis}\\
    \mathcal{B}(x)^{(l)}_{m,d}
    &=\sum_{c}P^{(l)}_{dc}\,x^{(l)}_{m,c}
    +\delta_{l0}\delta_{m0}\,q_{d}
    +\sum_{k,a,b}B^{(l)}_{d,k,ab}\,z^{(l)}_{m,k,ab},
    \label{eq:sm-wbb-out}
\end{align}
with the per-degree matrices $A^{\mathrm{L}},A^{\mathrm{R}},P$, the bias $q$, and the mixing coefficients $B$ learnable.
A \emph{Wigner bilinear network} (WBN) is a composition $\pi\circ\mathcal{B}_{K}\circ\cdots\circ\mathcal{B}_{1}\circ\iota$ with equivariant linear maps $\iota:\mathcal{V}\to\mathcal{H}$ and $\pi:\mathcal{H}\to\mathcal{W}$.
Each block is equivariant: a rotation translates the argument of the synthesized fields, the point-wise product commutes with this translation, and the Haar projection~\eqref{eq:sm-wbb-analysis} intertwines it back to the coefficient action.

\begin{lemma}[Clebsch--Gordan trees span polynomial equivariants]
\label{lem:sm-cg-span}
Let $\mathcal{V},\mathcal{W}$ be finite direct sums of SO(3) irreps.
Every polynomial SO(3)-equivariant map $P:\mathcal{V}\to\mathcal{W}$ is a finite linear combination of \emph{tree maps}: compositions of equivariant linear leaf maps $\mathcal{V}\to V_{l}$, iterated binary couplings~\eqref{eq:sm-cg-map}, and a final equivariant linear map into $\mathcal{W}$, together with constant terms in the trivial isotypic part of $\mathcal{W}$.
\end{lemma}

\begin{proof}
Decomposing $P$ into homogeneous parts and comparing degrees in $P(tx)$ shows each part is equivariant; the degree-zero part is an invariant constant.
A homogeneous part of degree $d$ is, by polarization, the diagonal restriction of a symmetric $d$-linear equivariant map, hence of an element of $\mathrm{Hom}_{SO(3)}(\mathcal{V}^{\otimes d},\mathcal{W})$; it therefore suffices to span this space by trees, and, after decomposing $\mathcal{V}$ and $\mathcal{W}$ into irreps, to span $\mathrm{Hom}_{SO(3)}(V_{l_1}\otimes\cdots\otimes V_{l_d},V_L)$.
We induct on $d$.
For $d=1$ Schur's lemma leaves only multiples of the identity.
For $d>1$, decompose $X=V_{l_1}\otimes\cdots\otimes V_{l_{d-1}}$ into irreducible summands with equivariant projections $p_{J,r}:X\to V_J$ and inclusions $i_{J,r}$ (the equivariant maps extracting and identifying the $r$-th copy of $V_J$ in the decomposition of $X$ into irreducibles, with $\sum_{J,r}i_{J,r}p_{J,r}=\mathrm{id}_X$); any $T\in\mathrm{Hom}_{SO(3)}(X\otimes V_{l_d},V_L)$ writes as $T=\sum_{J,r}T\circ(i_{J,r}p_{J,r}\otimes\mathrm{id})$, and each factor $T\circ(i_{J,r}\otimes\mathrm{id})$ lies in $\mathrm{Hom}_{SO(3)}(V_J\otimes V_{l_d},V_L)$, which by the multiplicity-free Clebsch--Gordan decomposition is spanned by $C^{L}_{Jl_d}$ when the triple is admissible and vanishes otherwise.
The induction hypothesis expresses each $p_{J,r}$ by trees on the first $d-1$ factors; composing with $C^{L}_{Jl_d}$ yields trees on all $d$ factors.
Evaluating on the diagonal $(x,\ldots,x)$ recovers the homogeneous polynomial parts and proves the claim.
\end{proof}

\begin{theorem}[Completeness of the WBN]
\label{thm:sm-wbn-complete}
Let $\mathcal{V},\mathcal{W}$ be finite direct sums of SO(3) irreps and let $F:\mathcal{V}\to\mathcal{W}$ be continuous and SO(3)-equivariant.
For every compact $\Omega\subset\mathcal{V}$ and every $\varepsilon>0$ there exist a band limit $L_{\max}$, a width $C$, a depth $K$, and WBN parameters such that
\begin{equation}
    \sup_{x\in\Omega}
    \bigl\|\pi\circ\mathcal{B}_{K}\circ\cdots\circ\mathcal{B}_{1}\circ\iota(x)-F(x)\bigr\|
    <\varepsilon .
    \label{eq:sm-wbn-uat}
\end{equation}
Moreover, every \emph{polynomial} SO(3)-equivariant map is realized exactly, with zero error.
\end{theorem}

\begin{proof}
\emph{Step 1 (reduction to polynomial equivariants).}
Let $\widehat\Omega=SO(3)\,\Omega$, compact by compactness of the group.
By the Stone--Weierstrass theorem there is a polynomial map $R:\mathcal{V}\to\mathcal{W}$ with $\sup_{\widehat\Omega}\|R-F\|<\delta$.
Its Reynolds average $P(x)=\int_{SO(3)}g^{-1}R(gx)\,dg$ is polynomial and equivariant, and since $F(x)=\int g^{-1}F(gx)\,dg$,
\[
    \sup_{x\in\Omega}\|P(x)-F(x)\|
    \le \Bigl(\sup_{g}\|g^{-1}\|_{\mathcal{W}}\Bigr)\,\delta ,
\]
which is smaller than $\varepsilon$ for $\delta$ small.
It therefore suffices to realize polynomial equivariants exactly.

\emph{Step 2 (one block realizes one coupling).}
Suppose two hidden channels of the current feature hold covariants $T_1(x)\in V_{l_1}$ and $T_2(x)\in V_{l_2}$, and fix an admissible output degree $L$.
Since the coupling map~\eqref{eq:sm-cg-map} is nonzero, its coefficient array is not identically zero, so there exist frames $|k_1|\le l_1$, $|k_2|\le l_2$ with
\begin{equation}
    \langle l_1 k_1\, l_2 k_2|L\,\kappa\rangle\neq0,
    \qquad \kappa=k_1+k_2 .
    \label{eq:sm-frame-choice}
\end{equation}
Use the maps~\eqref{eq:sm-wbb-lin} to load the two channels into one product pair,
\[
    f^{\mathrm{L}}_{a}(g)=\sum_{m_1}(T_1)_{m_1}D^{(l_1)}_{m_1k_1}(g),
    \qquad
    f^{\mathrm{R}}_{b}(g)=\sum_{m_2}(T_2)_{m_2}D^{(l_2)}_{m_2k_2}(g),
\]
all other coefficients set to zero.
By the product identity~\eqref{eq:sm-wigner-product},
\[
    h_{ab}(g)
    =\sum_{L'}\langle l_1 k_1\, l_2 k_2|L'\kappa\rangle\,
    \bigl(C^{L'}_{l_1l_2}(T_1,T_2)\bigr)_{M}\,D^{(L')}_{M\kappa}(g),
\]
and the analysis~\eqref{eq:sm-wbb-analysis} with the orthogonality~\eqref{eq:sm-wigner-orth} extracts
\[
    z^{(L)}_{M,\kappa,ab}
    =\langle l_1 k_1\, l_2 k_2|L\,\kappa\rangle\,
    \bigl(C^{L}_{l_1l_2}(T_1,T_2)\bigr)_{M}.
\]
The nonzero scalar~\eqref{eq:sm-frame-choice} is absorbed into $B^{(L)}_{d,\kappa,ab}$ in~\eqref{eq:sm-wbb-out}; choosing $P^{(l)}$ to preserve the existing channels and one fresh output channel $d$ stores exactly the coupling $C^{L}_{l_1l_2}(T_1,T_2)$.

\emph{Step 3 (finite stacks realize all trees).}
Fix a polynomial equivariant $P$ and expand it by Lemma~\ref{lem:sm-cg-span} into finitely many tree maps.
Assign one hidden channel to every node of every tree; choose $L_{\max}$ no smaller than the largest degree appearing and $C$ large enough to hold all channels.
The embedding $\iota$ places the leaf maps.
Proceeding by induction on tree depth, each block writes all couplings of the next depth into fresh channels -- Step~2, applied with a separate product pair $(a,b)$ and output channel per node, the mixing coefficients $B$ set to ignore all other pairs -- while preserving the channels already filled.
After finitely many blocks all nodes are present; the readout $\pi$ forms the finite linear combination of tree outputs, and the bias $q$ supplies the constant terms.
The stack realizes $P$ exactly, and with Step~1 the theorem follows.
\end{proof}

\begin{remark}[Frame truncation and sphere projection]
\label{rem:sm-wbn-truncation}
The full frame range enters the proof only through the choice~\eqref{eq:sm-frame-choice}.
Truncating the frames to $|k|\le k_{\max}$ restricts that choice, and at $k_{\max}=0$ it forces $k_1=k_2=\kappa=0$: the surviving coefficient $\langle l_1 0\, l_2 0|L 0\rangle$ vanishes whenever $l_1+l_2+L$ is odd, so a sphere-restricted bilinear nonlinearity cannot realize any parity-odd coupling, however deep the stack -- the incompleteness discussed in main-text Section~\ref{sec:s2-method}.
Already at $k_{\max}=1$ nonzero parity-odd choices become available, e.g.\ $\langle 1\,1\,1\,{-1}|1\,0\rangle=1/\sqrt{2}$ for the cross product $1\otimes1\to1$; the hypothesis of Theorem~\ref{thm:sm-wbn-complete}, however, requires the full range.
\end{remark}

\FloatBarrier
\section{Training and systems methods}
\label{sec:sm-training-systems}

\subsection{Training objective}
\label{sec:sm-training-objective}

DPA4 is trained as a conservative interatomic potential.
The neural network predicts a scalar total energy, and forces are obtained by differentiating this energy with respect to atomic positions.
The reported training runs use MAE losses with vector-norm force residuals.
For a mini-batch of configurations
$b=1,\ldots,B$, with $N_b$ atoms in configuration $b$, the objective is
\begin{equation}
    \mathcal{L}
    =
    \lambda_E
    \frac{1}{B}\sum_{b=1}^{B}
    \frac{|E_{\Theta,b}-E_b|}{N_b}
    +
    \lambda_F
    \frac{1}{\sum_{b=1}^{B}
        N_b} \sum_{b=1}^{B}\sum_{i=1}^{N_b} \left\| \mathbf{F}_{\Theta,bi}-\mathbf{F}_{bi} \right\|_2 + \lambda_{\Pi} \frac{1}{B}\sum_{b=1}^{B} \frac{\left\|\Pi_{\Theta,b}-\Pi_b\right\|_1}{9N_b}.
    \label{eq:sm-loss}
\end{equation}
Here $E_{\Theta,b}$, $\mathbf{F}_{\Theta,bi}$ and $\Pi_{\Theta,b}$ denote DPA4 predictions, while $E_b$, $\mathbf{F}_{bi}$ and $\Pi_b$ denote reference DFT labels.
The force residual is treated as a three-dimensional vector for each atom: the Euclidean norm is taken before averaging over atoms.
Energy and virial residuals use the MAE form with per-atom normalization.
The benchmark-specific weights, batch sizes and training lengths are listed in Tables~\ref{tab:sm_main_ablation_config}--\ref{tab:sm_matbench_config}.

\subsection{Warmup--stable--decay learning-rate schedule}
\label{sec:sm-wsd-schedule}

For long training runs, DPA4 uses a warmup--stable--decay schedule in which the learning rate first increases linearly, remains constant for the main training phase and is annealed only near the end of the run~\cite{wen2024wsd}.
Let $T$ be the total number of optimization steps, $T_{\mathrm{w}}$ the warmup length, $T_{\mathrm{d}}$ the decay length and $T_{\mathrm{s}}=T-T_{\mathrm{w}}-T_{\mathrm{d}}$ the stable length.
The
schedule used in the reported DPA4 training runs is
\begin{equation}
    \alpha(t)=
    \begin{cases}
        \alpha_{\mathrm{w}}
        +(\alpha_0-\alpha_{\mathrm{w}})t/T_{\mathrm{w}},
         & 0\le t<T_{\mathrm{w}},                             \\[3pt]
        \alpha_0,
         & T_{\mathrm{w}}\le t<T_{\mathrm{w}}+T_{\mathrm{s}}, \\[3pt]
        \alpha_{\mathrm{d}}(\tau),
         & T_{\mathrm{w}}+T_{\mathrm{s}}\le t<T,              \\[3pt]
        \alpha_{\min},
         & t\ge T,
    \end{cases}
    \label{eq:sm-wsd-piecewise}
\end{equation}
where $\alpha_{\mathrm{w}}$ is the initial warmup learning rate,
$\alpha_0$ is the stable-phase learning rate,
$\alpha_{\min}$ is the final learning rate and
\begin{equation}
    \tau
    =
    \operatorname{clip}
    \left(
    \frac{t-T_{\mathrm{w}}-T_{\mathrm{s}}}{T_{\mathrm{d}}},
    0,1
    \right).
    \label{eq:sm-wsd-tau}
\end{equation}
with cosine annealing in the decay phase,
\begin{equation}
    \alpha_{\mathrm{d}}(\tau)
    =
    \alpha_{\min}
    +
    \frac{\alpha_0-\alpha_{\min}}{2}
    \left(1+\cos \pi\tau\right).
    \label{eq:sm-wsd-cosine}
\end{equation}
Thus the stable phase carries most optimization steps, whereas the final cosine decay suppresses high-learning-rate oscillations before checkpoint selection.

\subsection{HybridMuon optimizer}
\label{sec:sm-hybrid-muon}

DPA4 is optimized with HybridMuon, a matrix-aware hybrid optimizer adapted from Muon~\cite{jordan2024muon} and scalable-Muon training studies~\cite{liu2025muon}.
The design separates two classes of trainable parameters.
Matrix-valued hidden transformations are routed to Muon, whereas biases, normalization scales, one-dimensional parameters and explicitly marked auxiliary parameters are routed to Adam or decoupled-weight-decay Adam.
This static routing is built on the first optimizer step and remains fixed during training, so the update rule for each parameter is independent of the current gradient value.

\paragraph{Matrix views and slice mode.}
Let a trainable tensor have an effective shape obtained by removing singleton dimensions.
HybridMuon interprets a rank-two effective shape as one matrix.
For higher-rank equivariant tensors, DPA4 uses slice mode: the leading dimensions index independent blocks, and Muon is applied separately to every trailing $(m,n)$ matrix.
Thus, for an effective shape $(b_1,\ldots,b_q,m,n)$, the
optimizer constructs
\begin{equation}
    B=\prod_{a=1}^{q}b_a \quad\text{independent matrix blocks}\quad G_t^{(b)}\in\mathbb{R}^{m\times n}, \qquad b=1,\ldots,B.
    \label{eq:sm-muon-slice-view}
\end{equation}
This choice is important for SO(2)-structured equivariant weights: degree and order-indexed blocks are updated independently, rather than being flattened into one large matrix that would mix unrelated representation strata.
In the degree-wise SO(3) channel maps used by the SO(2) pre- and post-projections and by the equivariant FFN, the weight tensor has shape $(L+1,C_{\mathrm{in}},F C_{\mathrm{out}})$; slice mode therefore applies Muon separately to the $(C_{\mathrm{in}},F C_{\mathrm{out}})$ matrix of each degree $l$.
Local SO(2) linear maps keep separate matrix parameters for the $m=0$ block and the constrained $|m|>0$ blocks, so the optimizer acts on these structured matrix blocks without collapsing distinct equivariant subspaces.

\paragraph{Muon update.}
For each Muon-routed block, the optimizer maintains a momentum buffer and forms
a Nesterov-style update
\begin{align}
    M_t^{(b)} & = \beta M_{t-1}^{(b)}+(1-\beta)G_t^{(b)}, \\ U_t^{(b)} & = \beta M_t^{(b)}+(1-\beta)G_t^{(b)}.
    \label{eq:sm-muon-nesterov}
\end{align}
The matrix $U_t^{(b)}$ is then orthogonalized by a Newton--Schulz polar iteration.
For the standard square or batched path, the iteration starts from
$X_0=U_t^{(b)}/\|U_t^{(b)}\|_{\mathrm{F}}$ and applies
\begin{equation}
    X_{k+1} = a_k X_k + \left(b_k A_k+c_k A_k^2\right)X_k, \qquad A_k=X_kX_k^{\mathsf{T}}.
    \label{eq:sm-newton-schulz}
\end{equation}
DPA4 uses a two-stage schedule: eight fast iterations with $(a,b,c)=(3.4445,-4.7750,2.0315)$ followed by two Newton polishing iterations with $(a,b,c)=(2,-1.5,0.5)$.
The resulting polar factor is denoted $Q_t^{(b)}$.

\paragraph{Rectangular Gram path.}
For rectangular matrices, HybridMuon uses a compiled Gram Newton--Schulz path following the fast polar-decomposition formulation used for Muon~\cite{dao2026gramnewtonschulz}.
The matrix is oriented so that $m\le n$, normalized in single precision and iterated in half precision using a fixed Polar-Express coefficient schedule.
Since the iteration depends on $XX^{\mathsf{T}}\in\mathbb{R}^{m\times m}$, rectangular blocks with the same smaller dimension can be column-padded, concatenated and orthogonalized in a single grouped call.
Padding only the larger dimension preserves the Frobenius
norm and the Gram matrix,
\begin{equation}
    [X\;0][X\;0]^{\mathsf{T}}=XX^{\mathsf{T}},
    \label{eq:sm-gram-padding}
\end{equation}
so truncating the padded columns after the iteration exactly recovers the
unpadded result while reducing the number of small GPU launches.

\paragraph{Update-RMS matching and Adam-family path.}
In the default match-RMS mode, the Muon update for an $m\times n$ block is
scaled as
\begin{equation}
    \Delta W_t^{(b)}
    =
    -\alpha_t\,\gamma\,\sqrt{\max(m,n)}\,Q_t^{(b)},
    \qquad \gamma=0.18.
    \label{eq:sm-match-rms}
\end{equation}
This coefficient follows the update-RMS calibration used in scalable Muon training~\cite{liu2025muon}: it brings the per-element magnitude of the orthogonalized matrix update onto the same learning-rate scale as AdamW-like updates.
Parameters routed to Adam use first and second moments in single precision, with bias correction and $\epsilon=10^{-20}$.
Decoupled weight decay is applied to Muon-routed matrices and to decay-enabled Adam-routed tensors, but not to one-dimensional Adam-routed parameters.

\subsection{Magma-lite update damping}
\label{sec:sm-magma-lite}

Conservative energy-gradient training produces gradients that include mixed coordinate--parameter derivatives and can exhibit large block-to-block variation.
DPA4 therefore augments the Muon path with a deterministic Magma-lite damping rule, adapted from momentum-aligned gradient masking~\cite{joo2026magma}.
For each Muon block, let $G_t^{(b)}$ be the current gradient and $M_t^{(b)}$ the momentum buffer after the update in Eq.~\eqref{eq:sm-muon-nesterov}.
The block alignment is
\begin{equation}
    \chi_t^{(b)}
    =
    \frac{
        \left\langle M_t^{(b)},G_t^{(b)}\right\rangle_{\mathrm{F}}
    }{
        \|M_t^{(b)}\|_{\mathrm{F}}\,
        \|G_t^{(b)}\|_{\mathrm{F}}+\varepsilon
    },
    \qquad
    \chi_t^{(b)}\in[-1,1].
    \label{eq:sm-magma-cosine}
\end{equation}
The score is mapped through a temperature-scaled sigmoid and stretched to
$[0,1]$,
\begin{equation}
    r_t^{(b)}
    =
    \operatorname{clip}
    \left[
        \frac{
            \sigma(\chi_t^{(b)}/\tau)-\sigma(-1/\tau)
        }{
            \sigma(1/\tau)-\sigma(-1/\tau)
        },
        0,1
        \right],
    \qquad \tau=2.
    \label{eq:sm-magma-score}
\end{equation}
An exponential moving average gives
\begin{equation}
    u_t^{(b)}
    =
    \rho_{\mathrm{ema}} u_{t-1}^{(b)}
    +(1-\rho_{\mathrm{ema}})r_t^{(b)},
    \qquad \rho_{\mathrm{ema}}=0.9,
    \label{eq:sm-magma-ema}
\end{equation}
and the final Muon update is rescaled by
\begin{equation}
    \Delta W_t^{(b)}
    \leftarrow
    \left[s_{\min}+(1-s_{\min})u_t^{(b)}\right]\Delta W_t^{(b)},
    \qquad s_{\min}=0.1.
    \label{eq:sm-magma-lite}
\end{equation}
This rule differs from stochastic masking: all blocks remain active, and poorly aligned updates are continuously damped rather than randomly skipped.
The nonzero lower bound prevents the optimizer from freezing a block entirely, which is important for MLIP training where force labels can make the alignment temporarily noisy without implying that the corresponding representation block should stop learning.

\subsection{Compiled conservative energy-gradient training}
\label{sec:sm-compiled-training}

The conservative energy-gradient path is compiled separately from direct-force or density-denoising modes.
Force matching requires differentiating through
\begin{equation}
    \mathbf{F}_{\Theta}
    =
    -\frac{\partial E_{\Theta}}{\partial R},
    \qquad
    \frac{\partial\mathcal{L}}{\partial\Theta}
    \supset
    \frac{\partial^2E_{\Theta}}{\partial R\,\partial\Theta}.
    \label{eq:sm-double-backward}
\end{equation}
A standard compiled forward/backward stack captures the forward graph and its first reverse-mode derivative, but it does not directly expose a nested coordinate derivative inside the compiled region.
DPA4 resolves this by first executing the energy-to-force derivative during symbolic tracing and then lowering the resulting tensor graph with PyTorch Inductor~\cite{ansel2024pytorch2}.

\paragraph{Tracing the conservative lower graph.}
The compiled function wraps the lower energy computation as a tensor-only map: from extended coordinates, atom types, neighbor indices and optional conditioning tensors to energies, forces and virials.
Before entering the traced function, the extended coordinates are rebound to a fresh leaf tensor.
This restart removes any upstream graph carried by data loading or neighbor construction while preserving a well-defined coordinate endpoint for $\partial E/\partial R$.
During training, the force construction keeps the coordinate-derivative graph alive; the outer loss backward pass can therefore differentiate the force residual into the model parameters, realizing the mixed derivative in Eq.~\eqref{eq:sm-double-backward}.

\paragraph{Symbolic tracing and higher-order differentiability.}
The implementation traces the conservative lower graph with symbolic shapes and real tensor inputs.
Real inputs are needed because compact edge construction contains data-dependent operations; after that control flow is resolved, the runtime dimensions are represented symbolically.
The trace uses a small five-frame representative batch chosen to avoid known symbolic-dimension collisions with singleton axes, charge--spin width, Cartesian coordinates and virial components.
The SiLU backward operation is decomposed into elementary pointwise operations before tracing, so the compiler receives an explicit first-derivative graph when the optimizer later requests a second derivative.

\paragraph{Preserving the force-loss gradient.}
When the coordinate derivative is traced in training mode, autograd inserts detach nodes around saved forward activations.
In the traced FX graph these nodes would become ordinary tensor operations and would sever the path from the force loss back to the parameters.
DPA4 removes only the detach nodes matching the saved-tensor topology and keeps user-intended detach operations unchanged.
The edited graph is then rebuilt into a fresh FX graph before compilation, so all compiler passes see a consistent node topology.

\paragraph{Dynamic edge representation and cache structure.}
The padded DeePMD neighbor list is converted inside the graph into a compact edge list.
Edge vectors are formed by differentiable indexed selection from the extended coordinate tensor, and a single masked sentinel edge is appended to every batch.
The sentinel edge guarantees a nonempty edge tensor under symbolic shapes while contributing exactly zero to downstream reductions.
Compiled callables are cached by graph topology, including training versus evaluation, the presence of atomic virial outputs and coordinate-correction inputs.
This multi-slot cache avoids recompilation when training is periodically interrupted by validation, while keeping distinct output signatures in separate compiled graphs.

\paragraph{Inductor configuration.}
The traced graph is lowered with dynamic-shape compilation.
The compiled path uses deterministic compilation settings rather than autotuning, enables shape padding for fluctuating symbolic dimensions and disables compiler features that interfere with higher-order autograd metadata or produce unstable large fused reduction kernels for this higher-order graph.
Training and evaluation use separate fusion limits because the training graph contains the second-derivative branch whereas the evaluation graph does not.
This systems design preserves the scalar energy-to-force relation while allowing the conservative training path to remain inside compiled GPU code.



\FloatBarrier
\section{Ablation study}
\label{sec:sm-ablation}

This section provides the complete ablation evidence behind Section~\ref{sec:ablation}.
The first subsection reports the cumulative ablation behind main-text Fig.~\ref{fig:ablation}, and the following six give the full mechanism-level sweeps for graph compilation, attention aggregation, multi-focus design, the low-rank edge--node SO(2)-equivariant product, $\Sph$ activation and the Wigner bilinear network.
The remaining subsections report model-selection and robustness studies.
All ablations use the same WBM-subsampled evaluation protocol.
Relative efficiency metrics are normalized within each controlled group under matched H20 hardware, batch-size, data-loading and precision settings.
Unless otherwise noted, \textit{Train time (rel.)
} denotes relative training wall-clock time, and
\textit{Test time (rel.)} denotes the wall-clock time of a full test pass over the
WBM-subsampled test set in DeePMD-kit.
Boldface and underlining denote the best and second-best values only for metrics in which ranking is explicitly highlighted.
In the ablation tables, \cmark{} and \xmark{} denote an enabled and a disabled setting, respectively.
N/A denotes a parameter or option that is not applicable to the corresponding setting.

\subsection{Cumulative ablation}
\label{sec:sm-cumulative-ablation}

This ablation is the numerical basis of main-text Fig.~\ref{fig:ablation}.
Starting from the complete interaction block, the four components that can be switched off without changing the dimensions of the equivariant representation are removed one after another, and every configuration is retrained under the identical protocol given in the cumulative-ablation column of Table~\ref{tab:sm_main_ablation_config} (Table~\ref{tab:sm_cumulative_ablation}).
All five configurations use a single focus stream, so the multi-focus design contributes nothing to the reported differences; it is assessed separately at matched SO(2) width in Section~\ref{sec:sm-multifocus}.

Removing all four components raises the energy MAE from 31.007 to 39.671~meV/atom and the force MAE from 39.169 to 44.581~meV/\AA{}.
Conversely, the four components together raise the parameter count by 1.7$\times$, from 1.19~M to 2.05~M, and the wall-clock time of a full test pass by 14\%.
The energy error increases monotonically along the chain, whereas the force error does not: removing the message node WBN leaves it marginally lower, so the two WBN insertions are separated by the force metric only in the controlled comparison of Table~\ref{tab:sm_wbn_ablation}.

\begin{table}[ht]
    \centering
    \scriptsize
    \begin{threeparttable}
        \caption{Cumulative ablation of the DPA4 architectural components.\tnote{a}}
        \label{tab:sm_cumulative_ablation}
        \begin{tabular}{@{}ccccccccc@{}}
            \toprule
            \makecell{\textbf{Message}\\\textbf{node WBN}}       &
            \makecell{\textbf{WBN in}\\\textbf{FFN}}             &
            \makecell{\textbf{Edge--node}\\\textbf{product}}     &
            \makecell{\textbf{Attention}\\\textbf{aggregation}}  &
            \textbf{E MAE}$\downarrow$                           &
            \textbf{F MAE}$\downarrow$                           &
            \makecell{\textbf{Train time}\\\textbf{(rel.)}$\downarrow$} &
            \makecell{\textbf{Test time}\\\textbf{(rel.)}$\downarrow$}  &
            \textbf{Params}                                                                        \\
            \midrule
            \cmark & \cmark & \cmark & \cmark & \textbf{31.007}    & \underline{39.169} & 1.00 & 1.00 & 2.05M \\
            \xmark & \cmark & \cmark & \cmark & \underline{31.585} & \textbf{38.957}    & 0.93 & 0.98 & 1.89M \\
            \xmark & \xmark & \cmark & \cmark & 33.226             & 39.918             & 0.87 & 0.93 & 1.22M \\
            \xmark & \xmark & \xmark & \cmark & 34.529             & 42.004             & 0.80 & 0.90 & 1.21M \\
            \xmark & \xmark & \xmark & \xmark & 39.671             & 44.581             & 0.72 & 0.88 & 1.19M \\
            \bottomrule
        \end{tabular}
        \begin{tablenotes}
            \item[a] Train time (rel.) and Test time (rel.) are normalized to the complete
            architecture of the first row.
            A disabled edge--node product denotes the scalar-scaling form that applies only the
            $l=0$ radial profile to all angular channels, and a disabled attention aggregation
            denotes scatter-sum aggregation.
            All rows use a single focus stream.
        \end{tablenotes}
    \end{threeparttable}
\end{table}

\FloatBarrier

\subsection{Graph compilation and training precision}
\label{sec:sm-compile-ablation}

This ablation quantifies the systems-level benefit of graph compilation and reduced-precision tensor-core execution (Table~\ref{tab:dpa4_compile_ablation}).
Relative to the non-compiled FP32 baseline, bf16 AMP alone gives a 1.43$\times$ training speedup and reduces peak training memory by 59\%.
Graph compilation provides a larger and complementary gain: compiled FP32 training gives a 1.61$\times$ speedup, and enabling TF32 in this compiled path increases the speedup to 1.82$\times$.
The largest practical improvement is obtained by combining compilation with bf16 AMP, which gives a 3.1$\times$ speedup and reduces peak memory by 60\%.
Thus, the compiled mixed-precision path makes conservative energy-gradient training more than three times faster while using only about 40\% of the baseline peak GPU memory.

The corresponding accuracy changes are small compared with the efficiency gain.
Against the FP32 baseline (27.603~meV/atom and 34.246~meV/\AA{}), bf16 AMP changes the energy and force MAEs by 1.7\% and 2.0\%, respectively.
The compiled-bf16 setting, which gives the strongest speed and memory improvement, changes them by 2.7\% and 1.8\%; when TF32 is also enabled, the energy MAE changes by 2.8\% and the force MAE by only 0.1\%.
These differences are consistent with small numerical and stochastic variation rather than systematic accuracy degradation.
Although the random seed is fixed and \texttt{torch.compile} preserves the mathematical computation graph, compilation can change kernel fusion, hardware-specific kernel selection, reduction ordering, and tensor-core dispatch; bf16 AMP and TF32 also change the effective arithmetic used by eligible matrix operations.
DPA4 limits this sensitivity by keeping geometric preprocessing and normalization operations such as RMSNorm in FP32, while allowing large matrix operations to use bf16 or TF32 where appropriate.
We therefore enable \texttt{torch.compile}, bf16 AMP, and TF32 in the following ablations and benchmark experiments unless otherwise specified.

\begin{table}[ht]
    \centering
    \scriptsize
    \begin{threeparttable}
        \caption{Ablation of graph compilation and training precision.\tnote{a}}
        \label{tab:dpa4_compile_ablation}
        \begin{tabular}{@{}ccccccc@{}}
            \toprule
            \textbf{Compile}                       &
            \textbf{bf16 AMP}                      &
            \textbf{TF32}                          &
            \textbf{E MAE}$\downarrow$             &
            \textbf{F MAE}$\downarrow$             &
            \textbf{Train time (rel.)}$\downarrow$ &
            \makecell{\textbf{Peak train}\\\textbf{mem. (rel.)}$\downarrow$}                                         \\
            \midrule
            \xmark                                  & \xmark & N/A   & 27.603 & 34.246 & 1.00          & 1.00          \\
            \xmark                                  & \cmark  & N/A   & 28.081 & 34.936 & 0.70          & 0.41          \\
            \cmark                                   & \xmark & \xmark & 28.217 & 34.787 & 0.62          & 0.77          \\
            \cmark                                   & \xmark & \cmark  & 28.190 & 34.860 & 0.55          & 1.00          \\
            \cmark                                   & \cmark  & \xmark & 28.355 & 34.864 & \textbf{0.32} & \textbf{0.40} \\
            \cmark                                   & \cmark  & \cmark  & 28.364 & 34.275 & \textbf{0.32} & \textbf{0.40} \\
            \bottomrule
        \end{tabular}
        \begin{tablenotes}
            \item[a] Train time (rel.) and peak train memory are normalized to the non-compiled
            FP32 baseline measured on NVIDIA H20 hardware.
            Peak train memory denotes peak GPU memory during training.
            N/A indicates that TF32 tensor cores are not applicable for the non-compiled execution path.
        \end{tablenotes}
    \end{threeparttable}
\end{table}

\FloatBarrier
\subsection{Attention aggregation}
\label{sec:sm-attention-aggregation}

The attention ablation isolates the aggregation rule while holding the feature dimension and focus count fixed within each pair of rows (Table~\ref{tab:dpa4_attention_ablation}).
Replacing scatter-sum aggregation with attention-weighted sum consistently reduces both energy and force MAEs across the 64-channel, 96-channel, and 96-channel 2-focus settings.
The improvement is substantial and stable: energy MAE decreases by 8.3--9.3\%, and force MAE decreases by 4.7--6.8\% relative to the corresponding scatter-sum controls.
The gain is obtained with a small computational cost, increasing training time by only 5--6\% and leaving test time within 1--5\% of the control models.
This result supports the design choice of computing attention from rotationally invariant scalar channels: the model gains adaptive neighbor selection while preserving equivariance of the higher-order SO(2) features.

\begin{table}[ht]
    \centering
    \scriptsize
    \begin{threeparttable}
        \caption{Ablation of attention aggregation.\tnote{a}}
        \label{tab:dpa4_attention_ablation}
        \begin{tabular}{@{}ccccccc@{}}
            \toprule
            \textbf{Feature dim.}                  &
            \textbf{No. focuses}                   &
            \textbf{Aggregation}                   &
            \textbf{E MAE}$\downarrow$             &
            \textbf{F MAE}$\downarrow$             &
            \textbf{Train time (rel.)}$\downarrow$ &
            \textbf{Test time (rel.)}$\downarrow$                                                                                 \\
            \midrule
            64                                     & 1 & Scatter sum            & 30.691          & 40.072          & 1.00 & 1.00 \\
            64                                     & 1 & Attention-weighted sum & \textbf{27.839} & \textbf{38.184} & 1.05 & 1.01 \\
            \midrule
            96                                     & 1 & Scatter sum            & 30.068          & 38.407          & 1.00 & 1.00 \\
            96                                     & 1 & Attention-weighted sum & \textbf{27.567} & \textbf{36.127} & 1.05 & 1.05 \\
            \midrule
            96                                     & 2 & Scatter sum            & 30.935          & 36.639          & 1.00 & 1.00 \\
            96                                     & 2 & Attention-weighted sum & \textbf{28.083} & \textbf{34.158} & 1.06 & 1.02 \\
            \bottomrule
        \end{tabular}
        \begin{tablenotes}
            \item[a] Within each feature-dimension and focus-count setting, Train time (rel.) and
            Test time (rel.) are normalized to the corresponding scatter-sum aggregation.
        \end{tablenotes}
    \end{threeparttable}
\end{table}

\FloatBarrier
\subsection{Multi-focus design}
\label{sec:sm-multifocus}

The multi-focus comparison separates per-focus feature width from the number of parallel equivariant focus channels (Table~\ref{tab:dpa4_focus_ablation}).
Increasing the width of a single focus stream rapidly increases parameter count and inference cost, but the corresponding accuracy gains are uneven.
By contrast, multi-focus variants increase the effective SO(2) convolution dimension through several narrower focus channels and often reach better accuracy--cost trade-offs.
At an SO(2) dimension of 192, the 96-channel 2-focus model gives the best overall result, improving energy MAE from 29.418 to 26.994~meV/atom and force MAE from 39.529 to 36.408~meV/\AA{} relative to the 64-channel 1-focus baseline.
It also outperforms the 192-channel 1-focus model with more than 56\% fewer trainable parameters, approximately 23\% lower training time, and approximately 34\% lower inference time.
All rows in this sweep use the same learning-rate setting; as larger-capacity variants often benefit from smaller tuned learning rates, the reported MAEs of the largest configurations may be mildly conservative.
Under a shared training recipe, however, the 96-channel 2-focus configuration gives the most balanced point in this sweep.
These trends are consistent with the intended role of focus competition, where parallel equivariant sub-channels specialize to different edge-local geometric motifs before the rotate-back step.

\begin{table}[ht]
    \centering
    \scriptsize
    \begin{threeparttable}
        \caption{Ablation of SO(2) feature width and focus count.\tnote{a}}
        \label{tab:dpa4_focus_ablation}
        \begin{tabular}{@{}cccccccc@{}}
            \toprule
            \textbf{Feature dim.}                  &
            \textbf{No. focuses}                   &
            \textbf{SO(2) dim.}                    &
            \textbf{E MAE}$\downarrow$             &
            \textbf{F MAE}$\downarrow$             &
            \textbf{Train time (rel.)}$\downarrow$ &
            \textbf{Test time (rel.)}$\downarrow$  &
            \textbf{Params}                                                                                                  \\
            \midrule
            64                                     & 1 & 64  & 29.418             & 39.529             & 1.00 & 1.00 & 1.9M  \\
            96                                     & 1 & 96  & 28.671             & 38.333             & 1.37 & 1.49 & 4.1M  \\
            64                                     & 2 & 128 & 28.126             & 38.123             & 1.55 & 1.55 & 3.2M  \\
            128                                    & 1 & 128 & 28.708             & 37.882             & 1.80 & 1.86 & 7.3M  \\
            64                                     & 3 & 192 & 27.890             & 37.283             & 1.95 & 2.14 & 4.5M  \\
            96                                     & 2 & 192 & \textbf{26.994}    & \textbf{36.408}    & 2.28 & 2.54 & 7.0M  \\
            192                                    & 1 & 192 & 27.286             & 36.477             & 2.98 & 3.86 & 16.0M \\
            64                                     & 4 & 256 & 27.821             & 37.605             & 2.43 & 2.65 & 5.7M  \\
            256                                    & 1 & 256 & 28.409             & 36.670             & 4.50 & 5.17 & 28.4M \\
            96                                     & 3 & 288 & \underline{27.168} & \underline{36.429} & 3.25 & 3.62 & 9.8M  \\
            \bottomrule
        \end{tabular}
        \begin{tablenotes}
            \item[a] Train time (rel.) and Test time (rel.) are normalized to the 64-channel
            1-focus baseline.
            Params denotes the number of trainable parameters.
            All rows use the same learning-rate setting.
        \end{tablenotes}
    \end{threeparttable}
\end{table}

\FloatBarrier
\subsection{Low-rank edge--node \texorpdfstring{SO(2)}{SO2}-equivariant product}
\label{sec:sm-radial-coupling}

This ablation tests the low-rank edge--node SO(2)-equivariant product (A1), namely how edge-side angular information conditions the node-side SO(2) message in the local frame (Table~\ref{tab:dpa4_radial_degree_ablation}).
The scalar-scaling baseline uses only the $l=0$ edge feature to scale edge messages across angular orders, whereas degree mixing builds a cross-degree kernel from the SO(2) Clebsch--Gordan coefficients and the $l>0$ edge spherical harmonics in the local frame, mixing input and output angular degrees at fixed $|m|$ before the SO(2) stack.
Making this kernel channel-dependent improves expressivity, but a per-channel dense kernel would inflate the parameter count by a factor of the hidden width $H$.
DPA4 instead parameterizes the kernel as a rank-$R$ factorization across the channel index, with $R$ scalar degree-pair coefficients contracted against a learnable channel basis of width $R\leq H$ (main-text Eq.~\ref{eq:dpa4-K-lowrank}).
Even the rank-$R=1$ form captures most of the benefit: it reduces energy MAE from 28.493 to 27.611~meV/atom and force MAE from 38.349 to 35.689~meV/\AA{} relative to scalar scaling, corresponding to 3.1\% and 6.9\% improvements, respectively (Table~\ref{tab:dpa4_radial_degree_ablation}).
This improvement is obtained at low additional cost, with training and test times increasing only to 1.12$\times$ and 1.10$\times$ the scalar-scaling baseline.
A compact low-rank edge--node product therefore provides a favorable accuracy--throughput trade-off before the kernel is made more expressive.

The remaining rank sweep shows that a more expressive edge--node product is not monotonically better.
Increasing the rank from 1 to 4 gives the best energy and force MAEs, but also raises the training cost to 1.86$\times$ the baseline; rank 8, rank 16, and the full kernel are still more expensive yet give worse force MAEs (Table~\ref{tab:dpa4_radial_degree_ablation}).
These results suggest that compact per-channel kernels act as a structural regularizer, providing enough channel-specific angular mixing without introducing many weakly constrained radial--angular interactions.

\begin{table}[ht]
    \centering
    \scriptsize
    \begin{threeparttable}
        \caption{Ablation of the low-rank edge--node SO(2)-equivariant product.\tnote{a}}
        \label{tab:dpa4_radial_degree_ablation}
        \begin{tabular*}{\linewidth}{@{\extracolsep{\fill}}ccccccc@{}}
            \toprule
            \textbf{Edge--node product} &
            \textbf{Per-channel ker.} &
            \textbf{Rank} &
            \textbf{E MAE}$\downarrow$ &
            \textbf{F MAE}$\downarrow$ &
            \textbf{Train time (rel.)}$\downarrow$ &
            \textbf{Test time (rel.)}$\downarrow$ \\
            \midrule
            Scalar scaling & \xmark & N/A & 28.493 & 38.349 & 1.00 & 1.00 \\
            Degree mixing & \xmark & N/A & 28.494 & 37.212 & 1.09 & 1.06 \\
            Degree mixing & \cmark & 1 & 27.611 & \underline{35.689} & 1.12 & 1.10 \\
            Degree mixing & \cmark & 2 & \underline{26.983} & 36.127 & 1.66 & 1.13 \\
            Degree mixing & \cmark & 4 & \textbf{26.556} & \textbf{35.675} & 1.86 & 1.14 \\
            Degree mixing & \cmark & 8 & 27.106 & 36.662 & 2.27 & 1.19 \\
            Degree mixing & \cmark & 16 & 26.984 & 36.401 & 3.09 & 1.34 \\
            Degree mixing & \cmark & Full & 28.011 & 37.400 & 3.36 & 2.22 \\
            \bottomrule
        \end{tabular*}
        \begin{tablenotes}
            \item[a] Train time (rel.) and Test time (rel.) are normalized to the scalar-scaling
            baseline.
            Scalar scaling uses only the $l=0$ edge feature, whereas degree mixing uses higher-degree edge-equivariant features to mix angular degrees in the local SO(2) frame.
        \end{tablenotes}
    \end{threeparttable}
\end{table}

\FloatBarrier

\subsection{\texorpdfstring{$\Sph$}{S2} activation and quadrature}

This ablation separates two coupled design choices in the spherical-grid nonlinearity: the model component to which $\Sph$ activation is applied and the quadrature rule used to project grid features back to equivariant coefficients (Table~\ref{tab:sm_quadrature_ablation}).
When $\Sph$ activation is restricted to the FFN component, replacing the latitude--longitude product grid with the Lebedev rule slightly reduces the energy MAE while leaving the force MAE and computational cost essentially unchanged.
Applying $\Sph$ activation additionally inside the SO(2) convolution component does not improve the overall accuracy--cost trade-off under this setting; with the product grid, both MAEs increase and the relative training and inference costs more than double.
In this expanded activation configuration, Lebedev quadrature reduces the energy error and lowers the relative training and inference costs compared with the product grid, although the force MAE remains higher than in the FFN-only configuration.
These results indicate that the FFN-only $\Sph$ activation configuration with Lebedev quadrature is both more accurate and cheaper than the SO(2)+FFN configurations tested here, and is therefore used in the released DPA4 variants.

\begin{table}[ht]
    \centering
    \scriptsize
    \begin{threeparttable}
        \caption{Ablation of $\Sph$ activation placement and quadrature rule.\tnote{a}}
        \label{tab:sm_quadrature_ablation}
        \begin{tabular}{@{}ccccccc@{}}
            \toprule
            \textbf{SO(2) S2 act.}                 &
            \textbf{FFN S2 act.}                   &
            \textbf{Quadrature}                    &
            \textbf{E MAE}$\downarrow$             &
            \textbf{F MAE}$\downarrow$             &
            \textbf{Train time (rel.)}$\downarrow$ &
            \textbf{Test time (rel.)}$\downarrow$                                                                           \\
            \midrule
            \xmark                                  & \cmark & Product & \underline{29.225} & \textbf{38.055}    & 1.00 & 1.00 \\
            \xmark                                  & \cmark & Lebedev & \textbf{28.992}    & \underline{38.103} & 0.98 & 1.01 \\
            \cmark                                   & \cmark & Product & 31.074             & 39.522             & 2.41 & 2.23 \\
            \cmark                                   & \cmark & Lebedev & 29.709             & 39.808             & 1.98 & 1.76 \\
            \bottomrule
        \end{tabular}
        \begin{tablenotes}
            \item[a] Train time (rel.) and Test time (rel.) are normalized to the FFN-only $\Sph$ row
            using the latitude--longitude product grid.
            The quadrature rule is used when projecting $\Sph$-grid features back to equivariant coefficients.
        \end{tablenotes}
    \end{threeparttable}
\end{table}

\FloatBarrier

\subsection{Wigner bilinear network}

This ablation isolates the two places where the Wigner bilinear network (WBN) enters an interaction block: the equivariant FFN, where the WBN grid replaces the $\Sph$-grid nonlinearity, and the post-aggregation message--node path, where a WBN grid couples the aggregated message with the destination node state (Table~\ref{tab:sm_wbn_ablation}).
Each of the two insertions is assessed with the other held fixed.

With the message node WBN disabled, the FFN nonlinearity produces the largest single effect: removing it raises the energy MAE by 2.77~meV/atom (7.4\%) and the force MAE by 1.51~meV/\AA{} (3.7\%) relative to the $\Sph$-grid FFN, and replacing that $\Sph$ grid by a WBN grid lowers both errors by a further 0.96~meV/atom (2.8\%) and 0.63~meV/\AA{} (1.6\%) for a 2\% increase in relative training time.
Enabling the message node WBN helps at either FFN setting.
On top of the $\Sph$-grid FFN it lowers the energy MAE by 0.79~meV/atom (2.3\%) and the force MAE by 1.10~meV/\AA{} (2.8\%), and on top of the WBN FFN it lowers the force MAE by a further 0.45~meV/\AA{} (1.2\%) while raising the energy MAE by 0.25~meV/atom (0.75\%).

The three configurations that contain a WBN grid therefore differ by at most 0.75\% in energy MAE and 1.2\% in force MAE, three to ten times less than the 2.3--7.4\% effects established above, so this experiment does not separate them.
Enabling both paths attains the lower force MAE of the two configurations that use a WBN grid in the FFN and is the setting used by the released DPA4 variants.
The cost of the added capacity falls almost entirely on the parameter count rather than on runtime: across the table the parameter count grows from 1.2~M to 2.0~M while the relative test time grows by only 4\%.

\begin{table}[ht]
    \centering
    \scriptsize
    \begin{threeparttable}
        \caption{Ablation of the Wigner bilinear network.\tnote{a}}
        \label{tab:sm_wbn_ablation}
        \begin{tabular}{@{}ccccccc@{}}
            \toprule
            \textbf{FFN nonlinearity}              &
            \textbf{Message node WBN}              &
            \textbf{E MAE}$\downarrow$             &
            \textbf{F MAE}$\downarrow$             &
            \textbf{Train time (rel.)}$\downarrow$ &
            \textbf{Test time (rel.)}$\downarrow$  &
            \textbf{Params}                                                                                  \\
            \midrule
            None        & \xmark & 37.213             & 40.653             & 1.00 & 1.00 & 1.2M \\
            $\Sph$ grid & \xmark & 34.448             & 39.142             & 1.14 & 1.01 & 1.4M \\
            WBN         & \xmark & \textbf{33.485}    & 38.517             & 1.16 & 1.02 & 1.9M \\
            $\Sph$ grid & \cmark  & \underline{33.662} & \textbf{38.044}    & 1.24 & 1.02 & 1.6M \\
            WBN         & \cmark  & 33.736             & \underline{38.069} & 1.23 & 1.04 & 2.0M \\
            \bottomrule
        \end{tabular}
        \begin{tablenotes}
            \item[a] Train time (rel.) and Test time (rel.) are normalized to the row without an FFN
            nonlinearity.
            The FFN nonlinearity column selects the grid on which the equivariant feed-forward
            nonlinearity is evaluated; when a WBN grid is used in the FFN it supersedes the
            $\Sph$-grid path.
        \end{tablenotes}
    \end{threeparttable}
\end{table}

\FloatBarrier

\subsection{Layers}

Interaction depth is the strongest determinant of accuracy among the structural hyperparameters (Tables~\ref{tab:sm_layers}--\ref{tab:sm_ffn_stack}).
Both errors fall most rapidly over the first few interaction blocks and then with diminishing returns, while training and inference cost grow approximately linearly with depth (Table~\ref{tab:sm_layers}); the depth of each DPA4 variant is therefore chosen to balance accuracy against cost rather than to minimize the error alone.
Within a block, additional SO(2) sublayers improve accuracy steadily at modest cost ($\sim$44\% training-time increase from 2 to 5 SO(2) sublayers; Table~\ref{tab:sm_so2_stack}), whereas deepening the feed-forward sublayer produces only small and partly non-monotonic changes (Table~\ref{tab:sm_ffn_stack}).
We accordingly retain a single feed-forward sublayer and tune the interaction depth and SO(2)-stack depth per variant.

\begin{table}[ht]
    \centering
    \scriptsize
    \caption{Ablation of interaction-block depth.}
    \label{tab:sm_layers}
    \begin{tabular*}{\linewidth}{@{\extracolsep{\fill}}ccccc@{}}
        \toprule
        \textbf{No. layers}               &
        \textbf{E MAE}$\downarrow$        &
        \textbf{F MAE}$\downarrow$        &
        \textbf{Train time (rel.)}$\downarrow$ &
        \textbf{Test time (rel.)}$\downarrow$                             \\
        \midrule
        1                                 & 37.125 & 45.552 & 1.00 & 1.00 \\
        2                                 & 29.816 & 39.968 & 1.23 & 1.36 \\
        3                                 & 28.036 & 37.946 & 1.70 & 1.71 \\
        4                                 & 27.207 & 36.628 & 2.18 & 2.13 \\
        5                                 & 27.419 & 35.853 & 2.55 & 2.66 \\
        6                                 & 27.087 & 35.603 & 2.89 & 3.02 \\
        7                                 & 27.019 & 35.281 & 3.43 & 3.47 \\
        8                                 & 26.549 & 34.628 & 3.80 & 3.84 \\
        9                                 & 26.818 & 34.523 & 4.14 & 4.41 \\
        10                                & 26.488 & 33.752 & 4.59 & 4.89 \\
        11                                & 26.552 & 33.731 & 5.07 & 5.29 \\
        12                                & 26.491 & 33.851 & 5.50 & 5.35 \\
        13                                & 26.488 & 34.282 & 5.91 & 6.20 \\
        14                                & 26.227 & 32.821 & 6.30 & 6.64 \\
        15                                & 26.868 & 33.224 & 6.66 & 7.75 \\
        \bottomrule
    \end{tabular*}
\end{table}

\begin{table}[ht]
    \centering
    \scriptsize
    \caption{Ablation of SO(2)-stack depth per interaction block.}
    \label{tab:sm_so2_stack}
    \begin{tabular*}{\linewidth}{@{\extracolsep{\fill}}ccccc@{}}
        \toprule
        \textbf{SO(2) layers}             &
        \textbf{E MAE}$\downarrow$        &
        \textbf{F MAE}$\downarrow$        &
        \textbf{Train time (rel.)}$\downarrow$ &
        \textbf{Test time (rel.)}$\downarrow$                             \\
        \midrule
        2                                 & 30.012 & 41.482 & 1.00 & 1.00 \\
        3                                 & 27.483 & 39.521 & 1.10 & 1.20 \\
        4                                 & 27.407 & 38.822 & 1.28 & 1.26 \\
        5                                 & 27.242 & 38.327 & 1.44 & 1.38 \\
        \bottomrule
    \end{tabular*}
\end{table}

\begin{table}[ht]
    \centering
    \scriptsize
    \caption{Ablation of FFN-stack depth per interaction block.}
    \label{tab:sm_ffn_stack}
    \begin{tabular*}{\linewidth}{@{\extracolsep{\fill}}ccccc@{}}
        \toprule
        \textbf{FFN layers}              &
        \textbf{E MAE}$\downarrow$       &
        \textbf{F MAE}$\downarrow$       &
        \textbf{Train time (rel.)}$\downarrow$ &
        \textbf{Test time (rel.)}$\downarrow$                             \\
        \midrule
        1                                & 29.527 & 37.408 & 1.00 & 1.00 \\
        2                                & 29.142 & 36.670 & 1.05 & 1.03 \\
        3                                & 28.889 & 36.627 & 1.13 & 1.05 \\
        4                                & 29.227 & 36.225 & 1.18 & 1.08 \\
        \bottomrule
    \end{tabular*}
\end{table}

\FloatBarrier
\subsection{Attention-aggregation design variants}

These sweeps vary the attention parameterization around the default single-head design, toggling the value projection, the output projection, the pre-mixing, and the number of heads for the 64-channel, 96-channel and 96-channel two-focus configurations (Tables~\ref{tab:sm_attention_64}--\ref{tab:sm_attention_96x2}).
The value projection is a learnable linear map applied to the message before the attention-weighted sum, the output projection is a channel mixing applied to the aggregated equivariant feature, and the pre-mixing is a cross-focus channel mixing applied to the input before computing attention.
Across all three configurations, enabling attention with one or more heads consistently outperforms the no-attention scatter-sum baseline, in agreement with the main attention ablation.
The minimal single-head form, without value, output or pre-mixing projections, attains the lowest or close to the lowest energy MAE; the additional projections and extra heads yield no consistent gain while increasing the parameter count.
Single-head attention without further projections is therefore retained as the default.

\begin{table}[ht]
    \centering
    \scriptsize
    \caption{Ablation of attention-aggregation variants for the 64-channel, 1-focus configuration.}
    \label{tab:sm_attention_64}
    \begin{tabular*}{\linewidth}{@{\extracolsep{\fill}}ccccccc@{}}
        \toprule
        \textbf{Attn. heads}       & \textbf{Attention}         & \textbf{Value proj.} &
        \textbf{Output proj.}      & \textbf{Pre mixing}        &
        \textbf{E MAE}$\downarrow$ & \textbf{F MAE}$\downarrow$                                                                            \\
        \midrule
        0                          & \xmark                      & N/A                  & N/A   & N/A   & 30.691          & 40.072          \\
        1                          & \cmark                       & \xmark                & \xmark & \xmark & \textbf{27.839} & 38.184          \\
        1                          & \cmark                       & \xmark                & \cmark  & \xmark & 29.214          & 38.250          \\
        1                          & \cmark                       & \cmark                 & \xmark & \xmark & 28.773          & 37.775          \\
        1                          & \cmark                       & \cmark                 & \cmark  & \xmark & 28.656          & 37.688          \\
        1                          & \cmark                       & \cmark                 & \cmark  & \cmark  & 28.985          & 37.242          \\
        2                          & \cmark                       & \xmark                & \xmark & \xmark & 28.661          & 37.896          \\
        2                          & \cmark                       & \xmark                & \cmark  & \xmark & \underline{28.634} & 38.035          \\
        2                          & \cmark                       & \cmark                 & \xmark & \xmark & 28.958          & \underline{37.056} \\
        2                          & \cmark                       & \cmark                 & \cmark  & \xmark & 28.936          & \textbf{36.939} \\
        2                          & \cmark                       & \cmark                 & \cmark  & \cmark  & 28.811          & 37.168          \\
        \bottomrule
    \end{tabular*}
\end{table}

\begin{table}[ht]
    \centering
    \scriptsize
    \caption{Ablation of attention-aggregation variants for the 96-channel, 1-focus configuration.}
    \label{tab:sm_attention_96}
    \begin{tabular*}{\linewidth}{@{\extracolsep{\fill}}ccccccc@{}}
        \toprule
        \textbf{Attn. heads}       & \textbf{Attention}         & \textbf{Value proj.} &
        \textbf{Output proj.}      & \textbf{Pre mixing}        &
        \textbf{E MAE}$\downarrow$ & \textbf{F MAE}$\downarrow$                                                                            \\
        \midrule
        0                          & \xmark                      & N/A                  & N/A   & N/A   & 30.068          & 38.407          \\
        1                          & \cmark                       & \xmark                & \xmark & \xmark & \textbf{27.567} & 36.127          \\
        1                          & \cmark                       & \xmark                & \cmark  & \xmark & 28.301          & 36.065          \\
        1                          & \cmark                       & \cmark                 & \xmark & \xmark & 28.578          & 36.350          \\
        1                          & \cmark                       & \cmark                 & \cmark  & \xmark & 28.430          & 36.321          \\
        1                          & \cmark                       & \cmark                 & \cmark  & \cmark  & 28.975          & 36.225          \\
        2                          & \cmark                       & \xmark                & \xmark & \xmark & 27.732          & 36.680          \\
        2                          & \cmark                       & \xmark                & \cmark  & \xmark & 28.074          & 36.288          \\
        2                          & \cmark                       & \cmark                 & \xmark & \xmark & \underline{27.571} & 36.119          \\
        2                          & \cmark                       & \cmark                 & \cmark  & \xmark & 27.968          & \textbf{35.457} \\
        2                          & \cmark                       & \cmark                 & \cmark  & \cmark  & 28.473          & 35.858          \\
        3                          & \cmark                       & \xmark                & \xmark & \xmark & 28.170          & 36.419          \\
        3                          & \cmark                       & \xmark                & \cmark  & \xmark & 28.016          & 36.017          \\
        3                          & \cmark                       & \cmark                 & \xmark & \xmark & 27.593          & \underline{35.639} \\
        3                          & \cmark                       & \cmark                 & \cmark  & \xmark & 27.722          & 35.809          \\
        3                          & \cmark                       & \cmark                 & \cmark  & \cmark  & 28.126          & 36.100          \\
        \bottomrule
    \end{tabular*}
\end{table}

\begin{table}[ht]
    \centering
    \scriptsize
    \caption{Ablation of attention-aggregation variants for the 96-channel, 2-focus configuration.}
    \label{tab:sm_attention_96x2}
    \begin{tabular*}{\linewidth}{@{\extracolsep{\fill}}ccccccc@{}}
        \toprule
        \textbf{Attn. heads}       & \textbf{Attention}         & \textbf{Value proj.} &
        \textbf{Output proj.}      & \textbf{Pre mixing}        &
        \textbf{E MAE}$\downarrow$ & \textbf{F MAE}$\downarrow$                                                                            \\
        \midrule
        0                          & \xmark                      & N/A                  & N/A   & N/A   & 30.935          & 36.639          \\
        1                          & \cmark                       & \xmark                & \xmark & \xmark & \underline{28.083} & 34.158          \\
        1                          & \cmark                       & \xmark                & \cmark  & \xmark & 28.449          & 34.508          \\
        1                          & \cmark                       & \cmark                 & \xmark & \xmark & 28.363          & 34.460          \\
        1                          & \cmark                       & \cmark                 & \cmark  & \xmark & 28.212          & 35.045          \\
        1                          & \cmark                       & \cmark                 & \cmark  & \cmark  & 28.885          & 35.387          \\
        2                          & \cmark                       & \xmark                & \xmark & \xmark & 28.335          & \underline{34.013} \\
        2                          & \cmark                       & \xmark                & \cmark  & \xmark & 28.401          & 34.361          \\
        2                          & \cmark                       & \cmark                 & \xmark & \xmark & \textbf{27.814} & \textbf{33.786} \\
        2                          & \cmark                       & \cmark                 & \cmark  & \xmark & 28.440          & 34.440          \\
        2                          & \cmark                       & \cmark                 & \cmark  & \cmark  & 28.968          & 35.507          \\
        3                          & \cmark                       & \xmark                & \xmark & \xmark & 28.506          & 34.560          \\
        3                          & \cmark                       & \xmark                & \cmark  & \xmark & 29.274          & 34.624          \\
        3                          & \cmark                       & \cmark                 & \xmark & \xmark & 28.907          & 34.412          \\
        3                          & \cmark                       & \cmark                 & \cmark  & \xmark & 28.609          & 34.869          \\
        3                          & \cmark                       & \cmark                 & \cmark  & \cmark  & 28.957          & 34.719          \\
        4                          & \cmark                       & \cmark                 & \cmark  & \cmark  & 28.736          & 34.603          \\
        6                          & \cmark                       & \xmark                & \xmark & \xmark & 28.757          & 34.289          \\
        6                          & \cmark                       & \xmark                & \cmark  & \xmark & 28.968          & 34.778          \\
        6                          & \cmark                       & \cmark                 & \xmark & \xmark & 29.027          & 34.313          \\
        6                          & \cmark                       & \cmark                 & \cmark  & \xmark & 28.397          & 34.541          \\
        6                          & \cmark                       & \cmark                 & \cmark  & \cmark  & 29.483          & 34.889          \\
        \bottomrule
    \end{tabular*}
\end{table}

\FloatBarrier
\subsection{Normalization placement}

Normalization placement is examined by applying RMSNorm as pre-normalization, post-normalization or both within the SO(2) and feed-forward sublayers (Table~\ref{tab:sm_stability}).
A single normalization per sublayer is consistently more accurate than applying both, which adds layers without benefit.
SO(2) post-normalization combined with feed-forward pre-normalization yields the lowest energy MAE and close to the lowest force MAE, and this placement is used throughout.

\begin{table}[ht]
    \centering
    \scriptsize
    \caption{Ablation of SO(2) and FFN normalization placement.}
    \label{tab:sm_stability}
    \begin{tabular*}{\linewidth}{@{\extracolsep{\fill}}cccccc@{}}
        \toprule
        \makecell{\textbf{SO(2)}\\\textbf{pre-norm}}  &
        \makecell{\textbf{SO(2)}\\\textbf{post-norm}} &
        \makecell{\textbf{FFN}\\\textbf{pre-norm}}    &
        \makecell{\textbf{FFN}\\\textbf{post-norm}}   &
        \textbf{E MAE}$\downarrow$                    &
        \textbf{F MAE}$\downarrow$                                                                      \\
        \midrule
        \cmark                                               & \xmark & \cmark  & \xmark & \underline{28.493} & 37.683          \\
        \xmark                                              & \cmark  & \cmark  & \xmark & \textbf{28.013} & \underline{37.624} \\
        \cmark                                               & \xmark & \xmark & \cmark  & 30.237          & 38.608          \\
        \xmark                                              & \cmark  & \xmark & \cmark  & 29.321          & \textbf{37.466} \\
        \cmark                                               & \xmark & \cmark  & \cmark  & 29.322          & 38.548          \\
        \cmark                                               & \cmark  & \cmark  & \xmark & 28.987          & 38.250          \\
        \xmark                                              & \cmark  & \cmark  & \cmark  & 28.808          & 37.840          \\
        \cmark                                               & \cmark  & \xmark & \cmark  & 29.841          & 39.068          \\
        \cmark                                               & \cmark  & \cmark  & \cmark  & 30.215          & 39.208          \\
        \bottomrule
    \end{tabular*}
\end{table}

\FloatBarrier
\subsection{Learning-rate scheduler comparison}

Under otherwise identical settings, the warmup--stable--decay (WSD) schedule lowers both the energy and force MAEs relative to a cosine schedule (Table~\ref{tab:sm_training_recipe}).
This advantage is consistent with its extended high-learning-rate phase and short terminal decay, and the WSD schedule is used for the released DPA4 variants.

\begin{table}[ht]
    \centering
    \scriptsize
    \caption{Ablation of the learning-rate scheduler.}
    \label{tab:sm_training_recipe}
    \begin{tabular*}{\linewidth}{@{\extracolsep{\fill}}ccc@{}}
        \toprule
        \textbf{LR scheduler}          &
        \textbf{E MAE}$\downarrow$     &
        \textbf{F MAE}$\downarrow$                                                        \\
        \midrule
        Cosine                          & 28.638          & 38.607          \\
        WSD                             & \textbf{27.828} & \textbf{37.022} \\
        \bottomrule
    \end{tabular*}
\end{table}

\FloatBarrier
\section{Inference benchmark protocol}
\label{sec:sm-inference}

The throughput reported in Section~\ref{sec:efficiency} is measured on diamond supercells, which give a single-composition system whose size can be varied over more than two orders of magnitude while the local environment stays fixed; at the 6~\AA{} cutoff used by DPA4 every atom carries 158 neighbours, so the test is representative of a comparatively dense environment.
Each point averages 30 single-point evaluations after warm-up, and the sweep for each model terminates at the largest supercell that fits in the memory of one GPU.
MACE-Omat-Small and MACE-Omat-Medium are the \texttt{mace-omat-0-small} and \texttt{mace-omat-0-medium} checkpoints of the MACE-OMAT-0 release~\cite{batatia2025crosslearning}, obtained from \url{https://github.com/ACEsuit/mace-foundations}; the same checkpoints supply the MACE-Omat rows of Table~\ref{tab:omat24}.
The EquiformerV3 timing uses the MPtrj checkpoint of the $L_{\max}=4$ architecture, which carries the same parameter count and layer structure as the OMat24 model of Table~\ref{tab:omat24} and therefore the same inference cost.

All ASE inference benchmarks were run on the same NVIDIA H20 hardware and base system environment with CUDA 12.8, Ubuntu 20.04.6 LTS and GCC 9.4.0.
The benchmarks used the ASE calculator interface~\cite{ase-paper}; DPA4 calculators used compiled inference with Python 3.13.13, PyTorch 2.11.0+cu128, ASE 3.28.0 and NumPy 2.4.4.
The MACE baselines~\cite{batatia2022mace,barroso2024open,batatia2025crosslearning} used MACE 0.3.15 with Python 3.11.15 and PyTorch 2.10.0+cu128; the cuEq variants used NVIDIA cuEquivariance-accelerated equivariant kernels~\cite{nvidia_cuequivariance}.
EquiformerV3~\cite{liao2026equiformerv3} used the \texttt{atomicarchitects/equiformer\_v3} implementation at commit \texttt{a7300c5} with Python 3.11.15 and PyTorch 2.8.0+cu128.
All ASE baseline environments used ASE 3.28.0 and NumPy 2.4.4.

\section{Model and training configurations}

\subsection{Ablation model configurations}

Table~\ref{tab:sm_main_ablation_config} gives the configurations for the cumulative ablation of main-text Fig.~\ref{fig:ablation}, whose four controlled switches are the message node WBN, the WBN grid in the FFN, the edge--node product and the attention aggregation, and for the mechanism ablations: graph compilation, attention aggregation, multi-focus design, the low-rank edge--node SO(2)-equivariant product, $\Sph$ activation and the Wigner bilinear network.
Table~\ref{tab:sm_supplementary_ablation_config} gives those for the model-selection and robustness ablations.
Entries marked by ``--'' are the controlled variables within the corresponding experiment family; all other entries define the shared reference setting.
Horizontal rules separate related groups of hyperparameters without adding category labels to the table body.

\begin{table}[p]
    \centering
    \scriptsize
    \setlength{\tabcolsep}{2pt}
    \caption{Ablation model configurations (1).}
    \label{tab:sm_main_ablation_config}
    \begin{threeparttable}
        \begin{tabular*}{\linewidth}{@{\extracolsep{\fill}}L{2.4cm}C{1.52cm}C{1.52cm}C{1.52cm}C{1.52cm}C{1.52cm}C{1.52cm}C{1.52cm}@{}}
            \toprule
            \textbf{Hyperparameter}       & \makecell{\textbf{Cumulative}\\\textbf{ablation}} & \makecell{\textbf{Compile/}\\\textbf{precision}} & \textbf{Attention}                                                           & \makecell{\textbf{Multi-focus}\\\textbf{SO(2)}} & \makecell{\textbf{Edge--node}\\\textbf{product}} & \makecell{\textbf{$\Sph$/}\\\textbf{quad.}} & \textbf{WBN} \\
            \midrule
            Feature dim.                  & 64                                                & 64                                               & 64/96/96                                                                     & --                                              & 64                                               & 64                                          & 64 \\
            No. focuses                   & 1                                                 & 1                                                & 1/1/2                                                                        & --                                              & 1                                                & 1                                           & 1 \\
            No. layers                    & 2                                                 & 3                                                & 3                                                                            & 2                                               & 3                                                & 3                                           & 2 \\
            SO(2) layers                  & 4                                                 & 4                                                & 4                                                                            & 4                                               & 4                                                & 4                                           & 4 \\
            FFN layers                    & 1                                                 & 1                                                & 1                                                                            & 1                                               & 1                                                & 1                                           & 1 \\
            \midrule
            Radial basis                  & Bessel                                            & Bessel                                           & Bessel                                                                       & Bessel                                          & Bessel                                           & Bessel                                      & Bessel \\
            No. radial bases              & 16                                                & 16                                               & 16                                                                           & 16                                              & 16                                               & 16                                          & 16 \\
            $L_{\max}$                    & 2                                                 & 3                                                & 3                                                                            & 3                                               & 3                                                & 3                                           & 2 \\
            $M_{\max}$                    & 1                                                 & 1                                                & 1                                                                            & 1                                               & 1                                                & 1                                           & 1 \\
            \midrule
            Edge--node product\tnote{g}   & --                                                & Degree mixing                                    & Scalar scaling                                                               & Scalar scaling                                  & --                                               & Degree mixing                               & Degree mixing \\
            Per-channel mod.              & --                                                & True                                             & N/A                                                                          & N/A                                             & --                                               & True                                        & True \\
            Rank                          & --                                                & 1                                                & N/A                                                                          & N/A                                             & --                                               & 1                                           & 2 \\
            Attn. heads                   & --                                                & 1                                                & --                                                                           & 1                                               & 1                                                & 1                                           & 1 \\
            Attention value/output proj.  & False                                             & False                                            & False                                                                        & False                                           & False                                            & False                                       & False \\
            Pre mixing                    & False                                             & False                                            & False                                                                        & False                                           & False                                            & False                                       & False \\
            Message node WBN              & --                                                & False                                            & False                                                                        & False                                           & False                                            & False                                       & -- \\
            S2 act.                       & --                                                & FFN only                                         & FFN only                                                                     & FFN only                                        & FFN only                                         & --                                          & -- \\
            WBN in FFN                    & --                                                & False                                            & False                                                                        & False                                           & False                                            & False                                       & -- \\
            Quadrature                    & Lebedev                                           & Lebedev                                          & Lebedev                                                                      & Product                                         & Lebedev                                          & --                                          & Lebedev \\
            Norm. placement\tnote{e}      & Post \& Pre                                       & Post \& Pre                                      & Post \& Pre                                                                  & Pre \& Pre                                      & Post \& Pre                                      & Pre \& Pre                                  & Post \& Pre \\
            Activation func.              & SiLU                                              & SiLU                                             & SiLU                                                                         & SiLU                                            & SiLU                                             & SiLU                                        & SiLU \\
            FFN hidden dim.               & Auto                                              & Auto\tnote{f}                                    & Auto                                                                         & Auto                                            & Auto                                             & Auto                                        & Auto \\
            SO(3) read-out                & False                                             & False                                            & False                                                                        & False                                           & False                                            & False                                       & False \\
            Output fitting dim.           & Auto                                              & Auto                                             & Auto                                                                         & Auto                                            & Auto                                             & Auto                                        & Auto \\
            Output fitting layers         & 1                                                 & 1                                                & 1                                                                            & 1                                               & 1                                                & 1                                           & 1 \\
            \midrule
            Compile                       & True                                              & --                                               & True                                                                         & True                                            & True                                             & True                                        & True \\
            bf16 AMP                      & True                                              & --                                               & True                                                                         & True                                            & True                                             & True                                        & True \\
            TF32 matmul                   & True                                              & --                                               & True                                                                         & True                                            & True                                             & True                                        & True \\
            LR scheduler                  & WSD                                               & Cosine                                           & Cosine                                                                       & WSD                                             & Cosine                                           & Cosine                                      & WSD \\
            Max. LR                       & $4\times10^{-4}$                                  & $4\times10^{-4}$                                 & \makecell{$4.5/4.2/3.5$\\$\times10^{-4}$}                                    & $4\times10^{-4}$                                & $4.5\times10^{-4}$                               & $4.5\times10^{-4}$                          & $4\times10^{-4}$ \\
            Min. LR                       & $1\times10^{-6}$                                  & $1\times10^{-6}$                                 & $1\times10^{-6}$                                                             & $1\times10^{-6}$                                & $1\times10^{-6}$                                 & $1\times10^{-6}$                            & $1\times10^{-6}$ \\
            Warmup ratio                  & 0.003                                             & 0.003                                            & 0.003                                                                        & 0.003                                           & 0.003                                            & 0.003                                       & 0.003 \\
            Decay ratio                   & 0.65                                              & N/A                                              & N/A                                                                          & 0.65                                            & N/A                                              & N/A                                         & 0.65 \\
            Decay type                    & Cosine                                            & N/A                                              & N/A                                                                          & Cosine                                          & N/A                                              & N/A                                         & Cosine \\
            Batch size (per GPU)\tnote{c} & $\lceil 900/N\rceil$                              & $\lceil 450/N\rceil$                             & \makecell{$\lceil 1000/N\rceil$\\$\lceil 700/N\rceil$\\$\lceil 400/N\rceil$} & $\lceil 600/N\rceil$                            & $\lceil 1000/N\rceil$                            & $\lceil 700/N\rceil$                        & $\lceil 900/N\rceil$ \\
            Training steps                & $1\times10^6$                                     & $1\times10^6$                                    & $1\times10^6$                                                                & $2\times10^6$                                   & $1\times10^6$                                    & $1\times10^6$                               & $5\times10^5$ \\
            No. GPUs                      & 1                                                 & 1                                                & 1                                                                            & 1                                               & 1                                                & 1                                           & 1 \\
            \midrule
            Loss                          & MAE                                               & MAE                                              & MAE                                                                          & MAE                                             & MAE                                              & MAE                                         & MAE \\
            Loss weights $(E,F,V)$        & 20, 20, 5                                         & 20, 20, 5                                        & 20, 20, 5                                                                    & 20, 20, 5                                       & 20, 20, 5                                        & 20, 20, 5                                   & 20, 20, 5 \\
            Optimizer                     & HybridMuon                                        & HybridMuon                                       & HybridMuon                                                                   & HybridMuon                                      & HybridMuon                                       & HybridMuon                                  & HybridMuon \\
            Muon mode                     & Slice                                             & Slice                                            & Slice                                                                        & Slice                                           & Slice                                            & Slice                                       & Slice \\
            Magma Lite                    & True                                              & True                                             & True                                                                         & True                                            & True                                             & True                                        & True \\
            Weight decay                  & $1\times10^{-3}$                                  & $1\times10^{-3}$                                 & $1\times10^{-3}$                                                             & $1\times10^{-3}$                                & $1\times10^{-3}$                                 & $1\times10^{-3}$                            & $1\times10^{-3}$ \\
            \midrule
            Cutoff (\AA)                  & 6                                                 & 6                                                & 6                                                                            & 6                                               & 6                                                & 6                                           & 6 \\
            \bottomrule
        \end{tabular*}
        \begin{tablenotes}
            \item[a] ``--'' indicates a controlled variable within the corresponding ablation
            family.
            \item[b] In the attention column, multi-value entries follow the 64-channel 1-focus,
            96-channel 1-focus, and 96-channel 2-focus settings.
            \item[c] $N$ denotes the number of atoms in each system; $\lceil\cdot\rceil$ rounds
            up to the nearest integer.
            \item[d] N/A denotes a parameter not used by the corresponding setting.
            \item[e] In the normalization-placement entry, the first term refers to the SO(2)
            subblock and the second term to the FFN subblock.
            \item[f] Auto denotes a hidden dimension inferred from the feature dimension and
            rounded up to a multiple of 32: $(8/3)d_\mathrm{feat}$.
            \item[g] Scalar scaling uses only the $l=0$ edge feature, whereas degree mixing
            uses higher-degree edge-equivariant features to mix angular degrees in the local SO(2)
            frame.
        \end{tablenotes}
    \end{threeparttable}
\end{table}

\begin{table}[p]
    \centering
    \scriptsize
    \renewcommand{\arraystretch}{0.97}
    \setlength{\tabcolsep}{3pt}
    \caption{Ablation model configurations (2).}
    \label{tab:sm_supplementary_ablation_config}
    \begin{threeparttable}
        \begin{tabular*}{\linewidth}{@{\extracolsep{\fill}}
                L{2.1cm}C{1.7cm}C{1.7cm}C{1.7cm}C{1.7cm}C{1.7cm}C{1.7cm}@{}} \toprule \textbf{Hyperparameter} & \textbf{Layers} & \makecell{\textbf{SO(2)}\\\textbf{stack}} & \makecell{\textbf{FFN}\\\textbf{stack}} & \textbf{Attention} & \makecell{\textbf{Norm.
                }\\\textbf{placement}} &
            \makecell{\textbf{LR}\\\textbf{scheduler}} \\
            \midrule
            Feature dim. & 64 & 64 & 64 & 64/96/96 & 64 & 64 \\
            No. focuses & 1 & 1 & 1 & 1/1/2 & 1 & 1 \\
            No. layers & -- & 2 & 2 & 3 & 3 & 3 \\
            SO(2) layers & 4 & -- & 3 & 4 & 4 & 4 \\
            FFN layers & 1 & 1 & -- & 1 & 1 & 1 \\
            \midrule
            Radial basis & Bessel & Bessel & Bessel & Bessel & Bessel & Bessel \\
            No. radial bases & 16 & 16 & 16 & 16 & 16 & 16 \\
            $L_{\max}$ & 3 & 3 & 3 & 3 & 3 & 3 \\
            $M_{\max}$ & 1 & 1 & 1 & 1 & 1 & 1 \\
            \midrule
            Edge--node product & Scalar scaling & Degree mixing & Degree mixing & Scalar scaling & Scalar scaling & Degree mixing \\
            Per-channel mod. & N/A & True & True & N/A & N/A & True \\
            Rank & N/A & 1 & 1 & N/A & N/A & 1 \\
            Attn. heads & 1 & 1 & 1 & -- & 1 & 1 \\
            Attention value/output proj. & False & False & False & -- & False & False \\
            Pre mixing & False & False & False & -- & False & False \\
            Message node WBN & False & False & False & False & False & False \\
            S2 act. & FFN only & FFN only & FFN only & FFN only & FFN only & FFN only \\
            WBN in FFN & False & False & False & False & False & False \\
            Quadrature & Product & Lebedev & Lebedev & Lebedev & Product & Lebedev \\
            Norm. placement & Pre \& Pre & Post \& Pre & Post \& Pre & Post \& Pre & -- & Post \& Pre \\
            Activation func. & SiLU & SiLU & SiLU & SiLU & SiLU & SiLU \\
            FFN hidden dim. & Auto & Auto & Auto & Auto & Auto & Auto \\
            SO(3) read-out & False & False & False & False & False & False \\
            Output fitting dim. & Auto & Auto & Auto & Auto & Auto & Auto \\
            Output fitting layers & 1 & 1 & 1 & 1 & 1 & 1 \\
            \midrule
            Compile & True & True & True & True & True & True \\
            bf16 AMP & True & True & True & True & True & True \\
            TF32 matmul & True & True & True & True & True & True \\
            LR scheduler & Cosine & WSD & WSD & Cosine & WSD & -- \\
            Max.
            LR & $5\times10^{-4}$ & $6.5\times10^{-4}$ & $4.5\times10^{-4}$ & \makecell{$4.5/4.2/3.5$\\$\times10^{-4}$} & $4\times10^{-4}$ & $4.5\times10^{-4}$ \\ Min.
            LR & $1\times10^{-6}$ & $1\times10^{-6}$ & $1\times10^{-6}$ & $1\times10^{-6}$ & $1\times10^{-6}$ & $1\times10^{-6}$ \\ Warmup ratio & 0.003 & 0.003 & 0.003 & 0.003 & 0.003 & 0.003 \\ Decay ratio & N/A & 0.65 & 0.65 & N/A & 0.65 & -- \\ Decay type & N/A & Cosine & Cosine & N/A & Cosine & -- \\ Batch size (per GPU) & $\lceil 512/N\rceil$ & $\lceil 2100/N\rceil$ & $\lceil 900/N\rceil$ & \makecell{$\lceil 1000/N\rceil$\\$\lceil 700/N\rceil$\\$\lceil 400/N\rceil$} & $\lceil 1000/N\rceil$ & $\lceil 1000/N\rceil$ \\ Training steps & $2\times10^6$ & $2\times10^6$ & $1\times10^6$ & $1\times10^6$ & $1\times10^6$ & $1\times10^6$ \\ No.
            GPUs & 1 & 1 & 1 & 1 & 1 & 1 \\ \midrule Loss & MAE & MAE & MAE & MAE & MAE & MAE \\ Loss weights $(E,F,V)$ & 20, 20, 5 & 20, 20, 5 & 20, 20, 5 & 20, 20, 5 & 20, 20, 5 & 20, 20, 5 \\ Optimizer & HybridMuon & HybridMuon & HybridMuon & HybridMuon & HybridMuon & HybridMuon \\ Muon mode & Slice & Slice & Slice & Slice & Slice & Slice \\ Magma Lite & True & True & True & True & True & True \\ Weight decay & $1\times10^{-3}$ & $1\times10^{-3}$ & $1\times10^{-3}$ & $1\times10^{-3}$ & $1\times10^{-3}$ & $1\times10^{-3}$ \\ \midrule Cutoff (\AA) & 6 & 6 & 6 & 6 & 6 & 6 \\
            \bottomrule
        \end{tabular*}
        \begin{tablenotes}
            \item[a] ``--'' indicates a controlled variable within the corresponding ablation
            family.
            \item[b] In the attention column, multi-value entries follow the 64-channel 1-focus,
            96-channel 1-focus, and 96-channel 2-focus settings.
            \item[c] $N$ denotes the number of atoms in each system; $\lceil\cdot\rceil$ rounds
            up to the nearest integer.
            \item[d] N/A denotes a parameter not used by the corresponding setting.
            \item[e] In the normalization-placement entry, the first term refers to the SO(2)
            subblock and the second term to the FFN subblock.
            \item[f] Auto denotes a hidden dimension inferred from the feature dimension and
            rounded up to a multiple of 32: $(8/3)d_\mathrm{feat}$.
            \item[g] Scalar scaling uses only the $l=0$ edge feature, whereas degree mixing
            uses higher-degree edge-equivariant features to mix angular degrees in the local SO(2)
            frame.
        \end{tablenotes}
    \end{threeparttable}
\end{table}

\subsection{Benchmark model configurations}

Table~\ref{tab:sm_model_config} lists the model and systems configurations of the four released DPA4 size classes, which are shared by every benchmark reported in this work; only the training schedule varies between datasets.
Tables~\ref{tab:sm_matbench_config}, \ref{tab:sm_omat24_config}, \ref{tab:sm_omol25_config} and \ref{tab:sm_spice_mace_off_config} give the training hyperparameters for the Matbench Discovery, OMat24, OMol25 and SPICE-MACE-OFF benchmarks, respectively.

\begin{table}[p]
    \centering
    \scriptsize
    \setlength{\tabcolsep}{3pt}
    \caption{DPA4 model and systems configurations of the four released size classes, shared by all benchmarks.}
    \label{tab:sm_model_config}
    \begin{threeparttable}
        \begin{tabular*}{\linewidth}{@{\extracolsep{\fill}}L{3.0cm}C{1.5cm}C{1.5cm}C{1.5cm}C{1.5cm}C{1.5cm}@{}}
            \toprule
            \textbf{Hyperparameter}      & \textbf{DPA4-Mini} & \textbf{DPA4-Neo}  & \textbf{DPA4-Air}  & \textbf{DPA4-Plus} & \textbf{DPA4-Pro}  \\
            \midrule
            Feature dim.                 & 32                 & 64                 & 64                 & 64                 & 128                \\
            No. focuses                  & 1                  & 2                  & 1                  & 1                  & 2                  \\
            No. layers                   & 2                  & 2                  & 3                  & 4                  & 6                  \\
            SO(2) layers                 & 3                  & 3                  & 4                  & 4                  & 4                  \\
            FFN layers                   & 1                  & 1                  & 1                  & 1                  & 1                  \\
            \midrule
            Radial basis                 & Bessel             & Bessel             & Bessel             & Bessel             & Bessel             \\
            No. radial bases             & 16                 & 16                 & 16                 & 16                 & 16                 \\
            $L_{\max}$                   & 2                  & 3                  & 3                  & 4                  & 5                  \\
            $M_{\max}$                   & 1                  & 1                  & 1                  & 1                  & 1                  \\
            \midrule
            Edge--node product\tnote{a}  & Degree mixing      & Degree mixing      & Degree mixing      & Degree mixing      & Degree mixing      \\
            Per-channel mod.             & True               & True               & True               & True               & True               \\
            Rank                         & 1                  & 1                  & 1                  & 2                  & 2                  \\
            Attn. heads                  & 1                  & 1                  & 1                  & 1                  & 1                  \\
            Attention value/output proj. & False              & False              & False              & False              & False              \\
            Pre mixing                   & False              & False              & False              & False              & False              \\
            Message node WBN    & True               & True               & True               & True               & True               \\
            WBN in FFN                   & True               & True               & True               & True               & True               \\
            Quadrature                   & Lebedev            & Lebedev            & Lebedev            & Lebedev            & Lebedev            \\
            Norm. placement\tnote{b}     & Post \& Pre        & Post \& Pre        & Post \& Pre        & Post \& Pre        & Post \& Pre        \\
            Activation func.             & SiLU               & SiLU               & SiLU               & SiLU               & SiLU               \\
            FFN hidden dim.              & Auto\tnote{c}      & Auto               & Auto               & Auto               & Auto               \\
            SO(3) read-out               & True               & True               & True               & True               & False              \\
            Output fitting dim.          & Auto               & Auto               & Auto               & Auto               & Auto               \\
            Output fitting layers        & 1                  & 1                  & 1                  & 1                  & 1                  \\
            \midrule
            Compile                      & True               & True               & True               & True               & True               \\
            bf16 AMP                     & True               & True               & True               & True               & True               \\
            TF32 matmul                  & True               & True               & True               & True               & True               \\
            \midrule
            Optimizer                    & HybridMuon         & HybridMuon         & HybridMuon         & HybridMuon         & HybridMuon         \\
            Muon mode                    & Slice              & Slice              & Slice              & Slice              & Slice              \\
            Magma Lite                   & True               & True               & True               & True               & True               \\
            Weight decay                 & $1\times10^{-3}$   & $1\times10^{-3}$   & $1\times10^{-3}$   & $1\times10^{-3}$   & $1\times10^{-3}$   \\
            \midrule
            Cutoff (\AA)                 & 6                  & 6                  & 6                  & 6                  & 6                  \\
            \midrule
            Params                       & 0.66M              & 1.1M               & 5.1M               & 8.8M               & 25.2M             \\
            \bottomrule
        \end{tabular*}
        \begin{tablenotes}
            \item[a] Scalar scaling uses only the $l=0$ edge feature, whereas degree mixing
            uses higher-degree edge-equivariant features to mix angular degrees in the local SO(2)
            frame.
            \item[b] In the normalization-placement entry, the first term refers to the SO(2)
            subblock and the second term to the FFN subblock.
            \item[c] Auto denotes a hidden dimension inferred from the feature dimension and
            rounded up to a multiple of 32: $(8/3)d_\mathrm{feat}$.
        \end{tablenotes}
    \end{threeparttable}
\end{table}

\begin{table}[p]
    \centering
    \scriptsize
    \setlength{\tabcolsep}{3pt}
    \caption{Matbench Discovery training hyperparameters.}
    \label{tab:sm_matbench_config}
    \begin{threeparttable}
        \begin{tabular*}{\linewidth}{@{\extracolsep{\fill}}L{2.5cm}C{1.55cm}C{1.55cm}C{1.55cm}C{1.55cm}@{}}
            \toprule
            \textbf{Hyperparameter}     & \textbf{DPA4-Neo}         & \textbf{DPA4-Air}         & \textbf{DPA4-Plus}        & \textbf{DPA4-Pro}        \\
            \midrule
            LR scheduler                & WSD                       & WSD                       & WSD                       & WSD                      \\
            Max. LR                     & $7\times10^{-4}$          & $4.2\times10^{-4}$        & $4.7\times10^{-4}$        & $4.3\times10^{-4}$       \\
            Min. LR                     & $1\times10^{-6}$          & $1\times10^{-6}$          & $1\times10^{-6}$          & $1\times10^{-6}$         \\
            Warmup ratio                & 0.003                     & 0.003                     & 0.003                     & 0.003                    \\
            Decay ratio                 & 0.65                      & 0.65                      & 0.65                      & 0.65                     \\
            Decay type                  & Cosine                    & Cosine                    & Cosine                    & Cosine                   \\
            Batch size (per GPU)\tnote{a} & $\lceil 3000/N\rceil$   & $\lceil 1800/N\rceil$     & $\lceil 1500/N\rceil$     & $\lceil 300/N\rceil$     \\
            Training epochs             & 100                       & 100                       & 100                       & 140                      \\
            No. GPUs                    & 2                         & 2                         & 4                         & 16                       \\
            \midrule
            Loss                        & MAE                       & MAE                       & MAE                       & MAE                      \\
            Loss weights $(E,F,V)$      & 20, 20, 5                 & 20, 20, 5                 & 20, 20, 5                 & 20, 20, 5                \\
            \bottomrule
        \end{tabular*}
        \begin{tablenotes}
            \item[a] $N$ denotes the number of atoms in each system; $\lceil\cdot\rceil$ rounds
            up to the nearest integer.
        \end{tablenotes}
    \end{threeparttable}
\end{table}

\begin{table}[p]
    \centering
    \scriptsize
    \setlength{\tabcolsep}{3pt}
    \caption{OMat24 training hyperparameters.}
    \label{tab:sm_omat24_config}
    \begin{threeparttable}
        \begin{tabular*}{\linewidth}{@{\extracolsep{\fill}}L{2.8cm}C{1.5cm}C{1.5cm}C{1.5cm}C{1.5cm}C{1.5cm}@{}}
            \toprule
            \textbf{Hyperparameter}       & \textbf{DPA4-Mini}     & \textbf{DPA4-Neo}     & \textbf{DPA4-Air}     & \textbf{DPA4-Plus}    & \textbf{DPA4-Pro}    \\
            \midrule
            LR scheduler                  & Cosine                 & Cosine                & Cosine                & Cosine                & Cosine               \\
            Max. LR                       & $1\times10^{-3}$       & $8\times10^{-4}$      & $4\times10^{-4}$      & $3\times10^{-4}$      & $2\times10^{-4}$     \\
            Min. LR                       & $1\times10^{-6}$       & $1\times10^{-6}$      & $1\times10^{-6}$      & $1\times10^{-6}$      & $1\times10^{-6}$     \\
            Warmup ratio                  & 0.003                  & 0.003                 & 0.003                 & 0.003                 & 0.003                \\
            Batch size (per GPU)\tnote{a} & $\lceil 10000/N\rceil$ & $\lceil 5000/N\rceil$ & $\lceil 3900/N\rceil$ & $\lceil 2000/N\rceil$ & $\lceil 400/N\rceil$ \\
            Training epochs               & 7                      & 7                     & 6                     & 5                     & 5                    \\
            No. GPUs                      & 2                      & 2                     & 4                     & 4                     & 16                   \\
            \midrule
            Loss                          & MAE                    & MAE                   & MAE                   & MAE                   & MAE                  \\
            Loss weights $(E,F,V)$        & 20, 20, 5              & 20, 20, 5             & 20, 20, 5             & 20, 20, 5             & 20, 20, 5            \\
            \bottomrule
        \end{tabular*}
        \begin{tablenotes}
            \item[a] $N$ denotes the number of atoms in each system; $\lceil\cdot\rceil$ rounds
            up to the nearest integer.
        \end{tablenotes}
    \end{threeparttable}
\end{table}

\begin{table}[p]
    \centering
    \scriptsize
    \setlength{\tabcolsep}{3pt}
    \caption{OMol25 training hyperparameters.}
    \label{tab:sm_omol25_config}
    \begin{threeparttable}
        \begin{tabular*}{\linewidth}{@{\extracolsep{\fill}}L{2.8cm}C{1.5cm}C{1.5cm}C{1.5cm}C{1.5cm}C{1.5cm}@{}}
            \toprule
            \textbf{Hyperparameter}       & \textbf{DPA4-Mini}     & \textbf{DPA4-Neo}     & \textbf{DPA4-Air}     & \textbf{DPA4-Plus}    & \textbf{DPA4-Pro}    \\
            \midrule
            LR scheduler                  & WSD                    & WSD                   & WSD                   & WSD                   & WSD                  \\
            Max. LR                       & $6\times10^{-4}$       & $8\times10^{-4}$      & $3.5\times10^{-4}$    & $3.5\times10^{-4}$    & $3\times10^{-4}$     \\
            Min. LR                       & $1\times10^{-6}$       & $1\times10^{-6}$      & $1\times10^{-6}$      & $1\times10^{-6}$      & $1\times10^{-6}$     \\
            Warmup ratio                  & 0.003                  & 0.003                 & 0.003                 & 0.003                 & 0.003                \\
            Decay ratio                   & 0.65                   & 0.65                  & 0.65                  & 0.65                  & 0.65                 \\
            Decay type                    & Cosine                 & Cosine                & Cosine                & Cosine                & Cosine               \\
            Batch size (per GPU)\tnote{a} & $\lceil 20000/N\rceil$ & $\lceil 6000/N\rceil$ & $\lceil 4000/N\rceil$ & $\lceil 2000/N\rceil$ & $\lceil 800/N\rceil$ \\
            Training epochs               & 18                     & 16                    & 14                    & 12                    & 10                   \\
            No. GPUs                      & 4                      & 8                     & 16                    & 24                    & 32                   \\
            \midrule
            Loss                          & MAE                    & MAE                   & MAE                   & MAE                   & MAE                  \\
            Loss weights $(E,F,V)$        & 10, 5, 0               & 10, 5, 0              & 10, 5, 0              & 10, 5, 0              & 10, 5, 0             \\
            \bottomrule
        \end{tabular*}
        \begin{tablenotes}
            \item[a] $N$ denotes the number of atoms in each system; $\lceil\cdot\rceil$ rounds
            up to the nearest integer.
        \end{tablenotes}
    \end{threeparttable}
\end{table}

\begin{table}[p]
    \centering
    \scriptsize
    \setlength{\tabcolsep}{3pt}
    \caption{SPICE-MACE-OFF training hyperparameters.}
    \label{tab:sm_spice_mace_off_config}
    \begin{threeparttable}
        \begin{tabular*}{\linewidth}{@{\extracolsep{\fill}}L{3.0cm}C{2.0cm}C{2.0cm}C{2.0cm}@{}}
            \toprule
            \textbf{Hyperparameter}       & \textbf{DPA4-Neo}     & \textbf{DPA4-Air}     & \textbf{DPA4-Plus}    \\
            \midrule
            LR scheduler                  & WSD                   & WSD                   & WSD                   \\
            Max. LR                       & $8\times10^{-4}$      & $5\times10^{-4}$      & $3.5\times10^{-4}$    \\
            Min. LR                       & $1\times10^{-6}$      & $1\times10^{-6}$      & $1\times10^{-6}$      \\
            Warmup ratio                  & 0.003                 & 0.003                 & 0.003                 \\
            Decay ratio                   & 0.65                  & 0.65                  & 0.65                  \\
            Decay type                    & Cosine                & Cosine                & Cosine                \\
            Batch size (per GPU)\tnote{a} & $\lceil 4000/N\rceil$ & $\lceil 2000/N\rceil$ & $\lceil 1000/N\rceil$ \\
            Training steps                & $2\times10^6$         & $2\times10^6$         & $2\times10^6$         \\
            No. GPUs                      & 1                     & 1                     & 1                     \\
            \midrule
            Loss                          & MAE                   & MAE                   & MAE                   \\
            Loss weights $(E,F,V)$        & 15, 20, 0             & 15, 20, 0             & 15, 20, 0             \\
            \bottomrule
        \end{tabular*}
        \begin{tablenotes}
            \item[a] $N$ denotes the number of atoms in each system; $\lceil\cdot\rceil$ rounds
            up to the nearest integer.
        \end{tablenotes}
    \end{threeparttable}
\end{table}
